\begin{dcomment}{exposition}{p57:On peut préciser la proposition 5.6 obtenue ainsi}{Small pieces and recognition of a class}
Grothendieck wants a supply of small generators which remains useful after one imposes several conditions on an object or on a morphism. The immediate example will be a class of cofibrations intersected with a class of weak equivalences. Before asking whether such an intersection generates a factorization system, he separates the more elementary question: can membership be recognized on sufficiently many small pieces? This is the purpose of the end of \S5. Its positive results survive the failure of the proposed factorization theorem in \S6.

Here and below all indexing categories are small in the fixed universe, and $\pi$ is infinite. Put $\mathcal K_\pi=\underline{\Filt}^{\pi}(\mathcal M)$ and write $\Ss_\pi(X)$ for the poset of subobjects of $X$ whose domains belong to $\mathcal K_\pi$. A poset is ``large with respect to $\pi$'' when every subset of cardinal at most $\pi$ has an upper bound; in current terminology it is $\pi^+$-directed. A $\pi$-adapted presentation has the additional continuity property
\[
 X_{\sup J}\simeq\varinjlim_{j\in J}X_j
 \quad\text{for every directed }J\text{ with }|J|\leq\pi,
 \qquad X_i\in\mathcal K_\pi.
\]
The last requirement is part of the definition. It will be precisely the requirement lost on manuscript page~83. The term ``$\pi$-accessible object'' means that its covariant hom functor preserves these large filtered colimits, hence means $\pi^+$-presentable in the modern cardinal convention. Accessibility of a \emph{class} in these pages is the stronger, explicitly embedded recognition condition $\mathrm b_\pi$: in every adapted presentation of $X$, the object $X$ belongs to $M$ exactly when the indices with $X_i\in M$ are cofinal. It must not be replaced without argument by the assertion that the full category on $M$ is accessible. The definition is recorded explicitly in Chapter~XIII, manuscript pages~86--89 \citep{Grothendieck1990XIII}.

The standing assumptions from the preceding pages have concrete roles. Cocompleteness supplies the colimits. The cardinal filtration makes each $\mathcal K_\pi$ essentially small and supplies presentations. Condition~(d) on manuscript pages~45--46 includes closure under colimits of diagrams with at most $\pi$ arrows, under subobjects, and under images of maps with small source. Condition~(e) says that filtered colimits in the arrow category preserve monomorphisms. Consequently the union of small subobjects is computed by their filtered colimit, and $\Ss_\pi(X)$ itself gives an adapted presentation. These hypotheses concern ordinary colimits, before homotopy enters the discussion.

There are two independent meanings of generation to keep in view. The $L$-envelope of $M_0$ is its closure under arbitrary small filtered colimits. The class $\widetilde C$ of arrows is a lifting closure, namely ${}^{\square}(C^{\square})$. A filtered presentation of a morphism need not be a sequence of pushout attachments. Nothing in Proposition~5.7 identifies these two closure operations.

The proposition compares three descriptions of an $L$-accessible class: an intrinsic recognition condition; closure of its bounded part under bounded filtered colimits; and a cofinality test on small subobjects. In condition~(c), $J$ is $\Ss_\pi(X)$. In condition~(c$'$), the phrase ``containing $M_0$'' is absent from the edition and has to be understood; otherwise the smallest filtered-colimit-closed class is the empty class. The same omission occurs in the parenthesis of Corollary~5.8. We retain these printed readings and record their intended meaning in the audit. The duplicate numbers 5.7 and 5.8 likewise belong to the edition: the proposition and corollary here are distinct from the lemmas with those numbers.
\end{dcomment}

\begin{dcomment}{exposition}{p60}{Why a closed cofinal set is the right index}
Grothendieck strengthens cofinality because it is not preserved by intersection alone. Two cofinal subsets of a directed set can be disjoint: the even and odd natural numbers already show this. His extra condition is continuity under bounded directed suprema. It ensures that alternating between two sets eventually produces a point belonging to both.

For clarity, the indexing hypotheses imply more than finite directedness. If $S\subset I$ has cardinal at most $\pi$, start with $S$, choose upper bounds for all pairs, repeat this operation countably many times, and call the union $K$. Since $\pi$ is infinite, $|K|\leq\pi$; the set $K$ is directed. Its supremum exists by $\pi$-adaptedness and bounds $S$. Thus $I$ is $\pi^+$-directed. If $J\subset I$ is cofinal and closed under the specified suprema, the same construction, with each new upper bound subsequently enlarged into $J$, shows that $J$ has the corresponding properties.

Now let $J,J'$ be closed and cofinal, and fix $i_0\in I$. Choose
\[
 i_0\leq j_0\leq j'_0\leq j_1\leq j'_1\leq\cdots,
 \qquad j_n\in J,\quad j'_n\in J'.
\]
Both subsequences have the same supremum. Closure places that supremum in both sets, and it still lies above $i_0$. This proves cofinality of $J\cap J'$; closure is immediate. The argument uses only a countable supremum, so no regularity assumption on $\pi$ is needed. The ambient convention $\pi\geq\aleph_0$ matters: the result, isolated from that convention and read for finite $\pi$, would fail for the even and odd subsets of $\mathbb N$.

Apply this observation to the small subobjects of an object $X$. If $N\subset\Filt^\pi$ is closed under filtered colimits of cardinal at most $\pi$, then
\[
 J_N(X)=\{Z\in\Ss_\pi(X):Z\in N\}
\]
is closed under the corresponding suprema. Therefore, whenever $J\subset\Ss_\pi(X)$ is closed and cofinal, $J_N(X)$ is cofinal if and only if $J\cap J_N(X)$ is cofinal. Membership is independent of the chosen continuous cofinal subsystem. This is the precise reason for the closure assumption on $M^\pi$ in the corrected version of Proposition~5.6.

The opening zigzag of these pages explains the remaining comparison with an arbitrary adapted presentation. At a cardinal for which the $X_i$ and their images are $\pi$-accessible, a map from $\operatorname{Im}(X_i\to X)$ into the colimit factors through a later stage. One must also enlarge that stage so that the composite from $X_i$ equals the transition map, rather than merely having the same composite to $X$. Preservation of the filtered colimit by $\operatorname{Hom}(X_i,-)$ supplies this equality. Repeating the construction gives the interlacing diagram printed on p.~34. At its countable supremum the two colimits agree; the image-side colimit maps monomorphically to $X$ by condition~(e). Hence there are cofinally many monomorphic stages. This fills the implicit equality step without changing Grothendieck's argument.
\end{dcomment}

\begin{dcomment}{exposition}{p64}{Integrating small presentations}
Grothendieck now wants to flatten a filtered diagram of objects which themselves admit small presentations. His category $\mathcal C$ remembers both the outer index and a small object mapping into that stage. Keeping these maps, rather than choosing presentations separately and hoping they are compatible, makes the construction functorial.

The notation in Lemma~5.10 needs a restriction. The subobjects which form a subcategory of $\underline M^\pi/Y$ are
\[
 \Ss_\pi^{M^\pi}(Y)=\{Z\hookrightarrow Y:Z\in M^\pi\},
\]
not all of $\Ss_\pi(Y)$. For example, in sets let $M$ be the class of nonempty sets and let $Y$ be nonempty. The empty subobject of $Y$ belongs to $\Ss_\pi(Y)$ but is not an object of $\underline M^\pi/Y$. The proof and the subsequent displayed formula use the restricted collection, so the repair is unambiguous.

Under the standing subobject and image hypotheses the intended lemma has a short proof. Any $Z\to Y$ with $Z\in M^\pi$ factors through its image, a small subobject by condition~(d). Since $Y\in M$, the recognition criterion enlarges this image to a subobject $T\hookrightarrow Y$ with $T\in M^\pi$. Such subobjects form a directed poset: first bound finitely many small subobjects in $\Ss_\pi(Y)$ and then enlarge the bound into the cofinal subset. Finally two maps $Z\rightrightarrows T$ over $Y$ are equal because $T\to Y$ is monic. Thus every finite diagram in the slice has a cone through this poset, and the inclusion is final. Its colimit is $Y$ because the small subobjects are cofinal among all small subobjects.

The formal lemma requested after Corollary~5.11 can also be made explicit. Let $p:E\to I$ be an opfibration, with $I$ filtered and every fibre nonempty and filtered. Given finitely many objects of $E$, move them by opcartesian arrows to one common fibre and take a cone there. Given parallel arrows of $E$, first equalize their images by moving farther in $I$, then use filteredness of the new fibre to equalize the resulting vertical arrows. Hence $E$ is filtered. For every $i\in I$, the category $(i\downarrow p)$ is nonempty, and the same argument makes it filtered; therefore $p$ is final. These are ordinary finality statements and do not require a homotopical version of Theorem~A.

For the category used here, objects are triples $(i,Z,u:Z\to X_i)$; a morphism to $(j,Z',u')$ consists of $a:i\to j$ and $v:Z\to Z'$ satisfying $u'v=X(a)u$. Postcomposition gives the opcartesian transport. Its tautological source functor is the dotted $\mathfrak z$ of the diagram. A cocone from this source functor to $T$ is exactly a compatible family of cocones on the fibres. Consequently
\[
 \varinjlim_{(i,Z,u)\in\mathcal C}Z
 \simeq\varinjlim_{i\in I}\left(\varinjlim_{Z\to X_i}Z\right)
 \simeq\varinjlim_{i\in I}X_i.
\]
This proves the interchange actually needed on page~65. Corollary~5.11 is printed with $X_i\in M^\pi$; the argument works for $X_i\in M$, which is the case required for Proposition~5.9. Its word ``fibrant'' should be ``cofibrant'', in the older terminology for an opfibration.
\end{dcomment}

\begin{dcomment}{exposition}{p65}{The cardinal regrouping in Proposition 5.9}
The last step of the proof turns an unrestricted filtered colimit into an adapted colimit whose stages remain bounded. This is a useful construction independently of the subobject description. Suppose
\[
 X=\varinjlim_{i\in I}Z_i,\qquad Z_i\in M^\pi,
\]
where $I$ is a directed poset. Let $P_\pi(I)$ consist of its nonempty directed subsets of cardinal at most $\pi$. Every subset of $I$ of cardinal at most $\pi$ lies in such a directed subset: adjoin chosen upper bounds of all pairs repeatedly and take the countable union. It follows that $P_\pi(I)$ is $\pi^+$-directed. The supremum of a directed family of at most $\pi$ members is its union, again a member of $P_\pi(I)$.

Set $Z_J=\varinjlim_{j\in J}Z_j$. Bounded filtered-colimit closure puts $Z_J$ in $M^\pi$. If $\mathcal J\subset P_\pi(I)$ is directed with $|\mathcal J|\leq\pi$, the universal property of a cocone gives
\[
 \varinjlim_{J\in\mathcal J}Z_J
 \simeq\varinjlim_{j\in\bigcup\mathcal J}Z_j
 =Z_{\bigcup\mathcal J}.
\]
Moreover $\varinjlim_{J\in P_\pi(I)}Z_J\simeq X$, because every object and every comparison in the original diagram occurs in some $J$. Thus the regrouped presentation is adapted and every one of its stages belongs to $M$. Condition $\mathrm b_\pi$ now implies $X\in M$. This supplies the continuity check implicit in the last displayed formula of page~65.

There is a genuine distinction between the hypotheses of Proposition~5.9 as printed and those used by its subsequent proof. The proposition lists only (a), (b), (c), and (d$_0$), whereas Lemma~5.10 is proved using images and the canonical subobject presentation from the stronger standing assumptions. The regrouping argument itself uses neither images nor exactness. Its preceding flattening step does. One complete version of the stated weakening is obtained by requiring that the objects of $\Filt^\pi$ be $\pi$-accessible. Then, for $Y\in M$, start with an adapted presentation of $Y$. The good indices are cofinal, and bounded filtered-colimit closure makes them closed under the required suprema. Maps from objects of $M^\pi$ into $Y$ factor through this subsystem, and equalities of maps hold at a sufficiently late stage, by presentability. These two properties show directly that $M^\pi/Y$ is filtered, that the chosen subsystem is final in it, and that its colimit is $Y$. The preceding block and the regrouping then prove Proposition~5.9 using (d$_0$) and this presentability assumption, without (d$_1$), (d$_2$), or (e).

Alternatively, under the full standing assumptions (a)--(e), the subobject proof in the preceding block works at the given cardinal. The assertion that both sets of additional hypotheses may be removed simultaneously is not established by the written sketch. This is a gap in its claimed generality, not a counterexample to its conclusion. Chapter~XIII, manuscript pages~90--97, explicitly retains the corresponding presentability condition when flattening slices \citep{Grothendieck1990XIII}. The operative form of Proposition~5.7 uses good cardinals, so its proof is unaffected by this distinction.

To finish that proposition, start with a class $M_0\subset\Filt^\pi$ and close it under bounded filtered colimits inside $\Filt^\pi$, obtaining $N$. Define $M_N$ by cofinality of $\Ss_\pi^N(X)$. The preceding arguments show that $M_N$ is $L$-stable and, at a good cardinal, $\pi$-accessible. It contains $M_0$, while each of its objects is a filtered colimit of objects of $N$, themselves in the $L$-envelope of $M_0$. Therefore $M_N$ is exactly that envelope. On a bounded object $X$, the poset $\Ss_\pi(X)$ has largest element $X$, so the recognition condition gives $M_N\cap\Filt^\pi=N$. This proves the inverse correspondence of Corollary~5.8, including its uniqueness assertion.
\end{dcomment}

\begin{dcomment}{exposition}{p67}{Two closure operations on accessible classes}
Grothendieck's letter $L$ records a deliberate strengthening of accessibility: the class is closed under \emph{all} small filtered colimits, including ordinary sequential colimits. Increasing the accessibility cardinal does not remove that requirement. Conversely, once a class is generated under filtered colimits by a small family, one may choose a cardinal bounding that family and then a good cardinal for the ambient filtration. Proposition~5.7 supplies the recognition condition at that cardinal. This is why the absolute notion in Definition~5.12 does not depend on the particular cardinal filtration.

The equivalence between a small generating family and a bounded essentially small family also has a precise size content. Each filtration stage is essentially small, so its members have a small set of representatives. Conversely, choose one representative of every isomorphism class in a given essentially small family. Each representative belongs to some stage, and the small set of those cardinal bounds has a common bound in the universe. Replacing the family by its isomorphism closure changes none of its colimits. These are the size reductions used in clauses (b) and (b$'$), rather than an assertion that a potentially large collection of literally equal objects is a set.

The finite-union statement may at first seem surprising, since diagrams can contain objects from different classes. Its explanation is entirely about the index. If a directed set is covered by finitely many subsets, one subset is cofinal. Otherwise choose, for each subset, a point above which it never occurs, and then bound those finitely many points. No subset can contain that upper bound, a contradiction. A cofinal subset of a directed poset is itself directed, so a filtered diagram whose objects lie in $M_1\cup\cdots\cup M_n$ has a cofinal subdiagram lying in one $M_r$. If every $M_r$ is $L$-stable, its colimit lies in that $M_r$ and hence in the union.

The same finite-colour argument proves the recognition property in an adapted presentation. If the union of the good-index subsets is cofinal, apply the lemma inside that union; one good-index subset is then cofinal in the union and hence in the original index. The individual good-index subsets need not all be filtered before one knows that they are cofinal. The parenthesis ``filtrantes'' on page~72 is therefore unnecessary and, if read as an unconditional assertion, incorrect. For example a diagram of sets indexed by the fork $a,b\leq t$, with two singleton stages mapping to the two different elements of a two-element terminal stage, has two singleton indices without a common singleton upper bound. The finite-cover lemma does not require that additional assertion.

For intersections, a different argument is needed. If $X$ belongs to every $M_\alpha$, each corresponding good-index subset is both cofinal and closed under bounded directed suprema, because $M_\alpha^\pi$ is closed under the matching colimits. The required common presentation is therefore obtained by intersecting closed cofinal subsets. This is the role of Lemma~5.15.
\end{dcomment}

\begin{dcomment}{exposition}{p71}{The transfinite intersection argument and its bound}
Grothendieck is arranging one approximation which satisfies every condition at once. The induction works because a limit of approximations eventually satisfying a fixed condition still satisfies that condition. Here is the proof with its terminal index included, so that the last intersection is obtained as well as all its proper initial intersections.

Write the indexing set as an ordinal $\lambda$ with $|\lambda|\leq\pi$, and put
\[
 K_\alpha=\bigcap_{\beta<\alpha}I_\beta\qquad(\alpha\leq\lambda).
\]
We prove by transfinite induction that every $K_\alpha$ is closed and cofinal. For $\alpha=0$ it is $I$. Closure holds for every intersection. At a successor ordinal it follows from the two-set lemma that $K_{\alpha+1}=K_\alpha\cap I_\alpha$ is cofinal. Suppose now that $\delta\leq\lambda$ is a nonzero limit and that all the earlier intersections are cofinal. Given $i\in I$, construct an increasing family $(x_\beta)_{\beta<\delta}$ with $x_0\geq i$ and $x_\beta\in K_\beta$. At a successor stage enlarge $x_\beta$ into $K_{\beta+1}$. At a limit stage $\eta<\delta$, take the supremum of the previous terms. This supremum exists since the chain has cardinal at most $\pi$. For each $\gamma<\eta$, its tail from $\gamma$ onward lies in $K_\gamma$; closure therefore places the supremum in every $K_\gamma$, hence in $K_\eta$.

Finally take $x=\sup_{\beta<\delta}x_\beta$. The same tail argument puts $x$ in $K_\gamma$ for every $\gamma<\delta$, and hence in $K_\delta$. It also lies above $i$. This proves the limit step and gives cofinality of $K_\lambda$, the desired full intersection. In the manuscript's version the line referring to $I_{\beta_0}$ is naturally read with $J_{\beta_0}$, as in the following line; alternatively one starts the tail after $\beta_0$. Neither adjustment changes the proof. What matters is to run the induction through the order type of the entire indexing set, not merely through its elements.

The cardinal restriction is exact at this level of generality. Let $T$ be a set of cardinal strictly larger than $\pi$, and let $I$ be the poset of subsets of $T$ of cardinal at most $\pi$. It is directed, and a directed family of at most $\pi$ such subsets has its union as supremum. For $t\in T$, the set
\[
 I_t=\{S\in I:t\in S\}
\]
is cofinal and closed under these suprema. Yet $\bigcap_{t\in T}I_t$ is empty, because no member of $I$ contains all of $T$. Thus one cannot intersect an arbitrarily large family at a fixed cardinal by this argument.

Corollary~5.16 is nevertheless an absolute small-intersection theorem. Given a small family of $L$-accessible classes, choose a common sufficiently large good cardinal witnessing their accessibility and bounding the size of the family. Apply Proposition~5.13 there. The witness cardinal may change; the absolute class does not. This is the flexibility that will allow several properties of arrows to be imposed simultaneously.
\end{dcomment}

\begin{dcomment}{examples}{p73}{Why finite unions are the natural statement}
Grothendieck's finite-union theorem has a useful algebraic test. Fix a prime $p$ and work in abelian groups with the cardinal filtration by underlying sets of cardinal at most $\pi$, where $\pi$ is infinite. For $n\geq1$, let
\[
 M_n=\{A:p^nA=0\}.
\]
The class is closed under filtered colimits: every element of a filtered colimit is represented by an element of a stage, and multiplication by $p^n$ is zero at that stage. It is also recognized by small subgroups. If $p^nA=0$, all subgroups satisfy the same equation. Conversely, if small subgroups annihilated by $p^n$ are cofinal, every element of $A$ belongs to one of them, so it is annihilated by $p^n$. These good subgroups are closed under bounded directed unions. At this cardinal the bounded groups are $\pi^+$-presentable: a group with at most $\pi$ elements has a presentation with at most $\pi$ generators and relations, and a map or an equality out of such a presentation is settled at one stage of a $\pi^+$-filtered diagram. Thus Proposition~5.7 applies and each $M_n$ is $L$-$\pi$-accessible.

The increasing union $\bigcup_nM_n$ is the class of groups annihilated by some finite power of $p$, with a bound allowed to depend on the whole group. It is not $L$-stable. Consider
\[
 \mathbb Z/p\longrightarrow\mathbb Z/p^2\longrightarrow
 \mathbb Z/p^3\longrightarrow\cdots,
 \qquad [a]\longmapsto[pa].
\]
All arrows are injective. Sending $[a]\in\mathbb Z/p^n$ to $a/p^n+\mathbb Z$ identifies its colimit with $\mathbb Z[1/p]/\mathbb Z$. For every $n$ this group has an element of order $p^{n+1}$; consequently no power $p^n$ kills the whole colimit. Every stage is in the union and the colimit is not. This computes the failure of an infinite-union analogue without any size pathology.

The example also explains the distinction between a global bound and a local condition. The class of \emph{all} $p$-primary torsion groups is $L$-stable: each element is annihilated by some power of $p$, and an element in a filtered colimit comes from a stage. Its small subgroups give the same cofinality test. Passing from bounded exponent to elementwise finite exponent changes the class, and is exactly what happens in the displayed colimit. Accessibility alone does not authorize exchanging these two quantifiers.
\end{dcomment}

\begin{dcomment}{context}{p74}{The passage to diagrams comes from SGA 4}
Grothendieck needs the preceding theory in an arrow category, since weak equivalences and cofibrations are classes of arrows. His citation is precise: SGA~4, Expos\'e~I, Corollary~9.24 gives the objectwise cardinal filtration on categories of sections, under the stated accessibility and filtered-colimit hypotheses, after the bound $\pi\geq2^{c_0}$ \citep[Expos\'e~I, 9.22--9.24]{SGA4I1972}. Ordinary functor categories are the case of a constant fibred category. The enlarged bound pays for simultaneously choosing the objects, the structure maps, and the equations making them a diagram; small approximations of individual values alone do not provide those equations.

The other transferred properties can be checked directly. Colimits in $\mathcal M^S$ are computed at each object of $S$. A sufficiently small colimit of diagrams whose values lie in $\Filt^\pi$ again has its values there. Essential smallness follows by choosing a small skeleton of $\Filt^\pi$: a functor assigns one of its objects to each object of $S$ and one of its arrows to each arrow of $S$, subject to equations. Since $S$ is small, the resulting set of choices is small. Images and subobjects are also computed pointwise under the standing hypotheses. For images, the arrow induced from $\operatorname{Im}(F(s)\to G(s))$ to $\operatorname{Im}(F(t)\to G(t))$ is forced by the image factorization and is unique after the monomorphic inclusion into $G(t)$; these unique arrows satisfy functoriality. This checks the part of condition~(d) that is easy to overlook.

A natural transformation is monic exactly when each component is monic in this setting. One direction follows directly from the definition. For the other, evaluation at $s$ has a left adjoint sending $A$ to the diagram $t\mapsto\coprod_{S(s,t)}A$; testing monicity against these diagrams detects monicity of the $s$-component. Filtered colimits therefore preserve monomorphisms objectwise whenever they do so in $\mathcal M$. The same pointwise argument transfers commutation of filtered colimits with finite limits when that stronger exactness condition is available.

For $S=[1]$, a bounded object of the arrow category is a map \emph{whose source and target are both bounded}. This concrete description should remain visible throughout \S6. A map from a small object into a fixed large object is not a bounded object of the arrow category. Nor does a family of such maps become a bounded presentation merely because it is continuous and its indexing poset has the correct suprema.

The relationship with modern accessible categories is easiest at good cardinals: the bounded diagram objects can then be compared with presentable objects, and the categorical smallness of all their structure maps is absorbed by increasing the regular cardinal. The historical filtration language retains this choice explicitly. It gives uniform control over several classes in the same ambient category, which is exactly the control required for $C\cap W$.
\end{dcomment}

\begin{dcomment}{exposition}{p77}{What the proposed generation theorem would have supplied}
Grothendieck wants the accessibility calculation to finish the existence problem for factorizations. If every arrow in $C$ lay in the lifting closure of a small bounded part $C_0$, the small object construction from \S1 would provide factorizations controlled by $C$. This is a substantially stronger assertion than saying that $C$ is a filtered-colimit closure of $C_0$.

Recall the order-reversing operations
\[
 \begin{aligned}
 C^{\square}&=\{p: p\text{ has the right lifting property against }C\},\\
 {}^{\square}R&=\{i:i\text{ has the left lifting property against }R\}.
 \end{aligned}
\]
In the manuscript $C_*=C^{\square}$ and $\widetilde C={}^{\square}(C^{\square})$. From $C_0\subset C\subset\widetilde {C_0}$ one obtains
\[
 C^{\square}\subset C_0^{\square}
       =(\widetilde {C_0})^{\square}\subset C^{\square}.
\]
Thus Corollary~6.2 is a correct formal consequence of the proposed inclusion. It supplies no independent evidence for that inclusion. The source itself marks the proof of Theorem~6.1 as withdrawn.

The distinction between lifting closure and closure under elementary constructions is also relevant to Corollary~6.3. Every left lifting class contains the isomorphisms and is closed under pushouts, composition, transfinite composition, and retracts. If a class is already closed under all filtered colimits, its closure under retracts can be proved without a separate hypothesis: an idempotent $e:f\to f$ in the arrow category has its image as the colimit of the sequence
\[
 f\xrightarrow e f\xrightarrow e f\xrightarrow e\cdots.
\]
A splitting of the idempotent identifies that image with the desired retract. Closure under filtered colimits and finite composition likewise supplies the limit stages of transfinite composition. These observations explain why the hypotheses in (b) look close to being sufficient. They do not produce a \emph{set} of lifting generators. The implication promised from (c$'$) to (a), and the assertion of a Quillen factorization obtained from \S1, cannot be certified by the cancelled proof of 6.1.

For comparison, modern ``accessible weak factorization system'' refers to accessibility of a chosen functorial factorization, not to an arbitrary class satisfying the manuscript's recognition condition. Rosick\'y's example of split monomorphisms and split epimorphisms in abelian groups gives an accessible weak factorization system which is not generated by a set \citep[Example~2.8]{Rosicky2017Accessible}. It is not by itself a counterexample to Theorem~6.1: split monomorphisms are not closed under all filtered colimits. An argument about this theorem must verify the additional $L$-condition and the embedded recognition condition.

The later counterexample discussed at manuscript page~84 does verify those conditions: pure monomorphisms in a finite presheaf diagram category provide the required separation between filtered generation and generation by lifting. Theorem~6.1 is therefore false in its proposed generality, rather than merely lacking the particular proof attempted here. The many valid cardinal and cofinality calculations of \S5 remain available; what fails is their conversion into uniformly small relative attachments.
\end{dcomment}

\begin{dcomment}{examples}{p81}{Exact filtered unions do not make every amalgam a subobject}
Grothendieck identifies two different requirements while trying to construct the next subobject of $B$. First, the approximating intersections $A\cap B_i$ must commute with bounded filtered unions. Second, after attaching $B_i$ to $A$ along that intersection, the pushout must still map monomorphically into $B$. The first requirement concerns a limit followed by a filtered colimit; the second concerns a pushout followed by an image comparison. They have different answers even in familiar algebraic categories.

Under condition~(e$'$), the identity
\[
 A\times_B\left(\varinjlim_iB_i\right)
       \simeq\varinjlim_i(A\times_BB_i)
\]
provides the desired continuity. Only this pullback along the fixed monomorphism $A\to B$ and filtered diagrams of subobjects of $B$ are needed. Commutation of filtered colimits with finite limits is sufficient. It also implies condition~(e): a map is monic exactly when its diagonal to its self-pullback is an isomorphism, and filtered colimits then preserve that characterization. The phrase ``dans $I$'' in the printed condition (e$_0$) should be read as ``dans $\mathcal M$''.

For the second requirement consider groups. Let $B=\mathbb Z\times\mathbb Z$, let $A=\mathbb Z\times0$, and let $B'=0\times\mathbb Z$. Both maps into $B$ are monomorphisms and $A\cap B'$ is the trivial group. Their pushout in groups is
\[
 A*B'\cong F(a,b),
\]
the free group on two generators. The canonical map to $B$ sends $a$ and $b$ to its two coordinate generators. It sends the commutator $aba^{-1}b^{-1}$ to zero, but that reduced word is nontrivial in the free group. Hence the canonical comparison $A*B'\to B$ is not monic. Meanwhile filtered colimits of groups commute with finite limits: both are formed on underlying sets, and a finite list of equations or a finite tuple of representatives is settled in one sufficiently late filtered stage. Thus even condition (e$_0$) does not repair this pushout problem.

In sets and presheaves of sets the corresponding pushout of two subobjects over their intersection is their union and does embed into $B$. The proof is componentwise: two elements in the disjoint union are identified precisely when they represent the same element in the intersection. This favourable geometry removes the obstruction just exhibited for groups. It still does not force the residual inclusion $A\cup B'\to B$ to remain in an arbitrarily chosen accessible class $C$. The pure-monomorphism counterexample later in the commentary occurs in such a favourable presheaf category, and therefore isolates that remaining obstruction.

A further hypothesis has slipped between the statement and this proof. The version of condition~(d) restated in Theorem~6.1 omits closure under subobjects, although page~80 uses it to conclude $A\cap B_i\in\Filt^\pi$. The original condition on manuscript pages~45--46 includes that closure. Retaining it makes this approximation step valid; it does not mend the final inference on page~84.
\end{dcomment}

\begin{dcomment}{development}{p83}{A continuous diagram outside the cardinal bound}
Grothendieck's footnote~28 locates the decisive obstruction accurately: the arrows $\gamma_i:C_i\to B$ are not known to be bounded objects of the arrow category. The following explicit test shows that the omitted bound cannot be recovered from continuity and an isomorphic colimit. We give the construction to isolate the inference, without making it serve as a counterexample to the entire theorem.

\textbf{An unbounded counterdiagram.}\ours
Let $\pi$ be an infinite cardinal, let $B$ be a set with $|B|>\pi$, and let $I=[B]^{\leq\pi}$, ordered by inclusion. In the arrow category of sets consider
\[
 \gamma_S:S\hookrightarrow B\qquad(S\in I),
\]
with transition squares given by inclusion on sources and the identity on $B$. The poset $I$ is $\pi$-adapted, the diagram preserves suprema of nonempty directed subsets of cardinal at most $\pi$, and its colimit is the identity of $B$. Nevertheless no $\gamma_S$ is an isomorphism.

\emph{Proof.} A family of at most $\pi$ subsets of cardinal at most $\pi$ has union of cardinal at most $\pi$. Thus $I$ is $\pi^+$-directed and directed unions of that size are its suprema. Colimits of the source sets in such a directed family are their unions. The target is the constant diagram at $B$, whose colimit is $B$ because the indexing category is nonempty and connected. Colimits in the arrow category are pointwise, proving the asserted continuity. Taking the union of all $S\in I$ gives $B$, and the colimit arrow is consequently $\operatorname{id}_B$. Finally every $S$ is a proper subset of $B$, so every stage fails to be an isomorphism.\hfill$\square$

For completeness, the class of isomorphisms has exactly the accessibility properties relevant here. It is closed under filtered colimits in the arrow category, since a compatible family of inverse maps induces the inverse of the colimit map. Its bounded part is closed under bounded filtered colimits. For an isomorphism $f:A\to D$, a small subarrow $A_0\to D_0$ can be enlarged to the isomorphism $f^{-1}(D_0)\to D_0$, still small; hence bounded isomorphic subarrows are cofinal. Conversely, if isomorphic subarrows are cofinal, their filtered colimit is $f$, so $f$ is an isomorphism. The standard cardinal filtration of the arrow category of sets has bounded objects $\pi^+$-presentable, and Proposition~5.7 therefore makes this class $L$-$\pi$-accessible in the manuscript's sense.

The displayed diagram does not contradict that recognition property: its target $B$ is too large, so none of its arrows belongs to $\Fl^\pi(\mathbf{Set})$. It does contradict the proposed inference after deleting precisely that hypothesis. Even the strongest possible conclusion about the colimit, namely that it is an identity, says nothing about cofinally many isomorphic stages in this situation.

This is an elementary diagnostic construction supplied here. We checked the bounded-presentation requirements in Chapter~XIII, manuscript pages~86--89, and the diagram-filtration statement of SGA~4 I~9.24; neither licenses the deleted hypothesis. We do not claim priority for this example. Its purpose is to make the cancelled inference visible before discussing a counterexample to the generation theorem itself.
\end{dcomment}

\begin{dcomment}{development}{p84:ce qui établit le lemme 6.4}{A presheaf counterexample to the abandoned generation theorem}
Grothendieck has isolated the precise defect of his argument, but he still asks whether the theorem might survive it. There is now a negative answer in a particularly elementary presheaf category. The modern obstruction is a theorem about pure monomorphisms. To apply it here, however, we must verify the manuscript's bounded cofinality condition; replacing that condition by the modern word ``accessible'' would leave the essential comparison undone.

Let $P$ be the three-object category with its two nonidentity arrows
\[
1\longleftarrow 0\longrightarrow 2,
\qquad \mathcal E=\mathrm{Set}^{P}.
\]
Thus an object is a span of sets $E_1\leftarrow E_0\to E_2$. This is the presheaf category on $P^{\mathrm{op}}$. Give $\mathcal E$ the filtration by the total cardinality of its three sets. For every infinite $\pi$, the objects of size at most $\pi$ are $\pi^+$-presentable: a natural transformation out of one of them is specified by at most $\pi$ elements and equations, all of which can be represented simultaneously at one stage of a $\pi^+$-filtered diagram. Subobjects, quotients, and colimits with at most $\pi$ arrows preserve this bound. All colimits and finite limits are computed on the three underlying sets. In particular, this example satisfies even the strengthened exactness assumptions introduced after Theorem~6.1.

An inclusion $U\hookrightarrow V$ in $\mathcal E$ is \emph{pure} if every finite system of equations in the three sorts, with parameters in $U$, that has a solution in $V$ already has a solution in $U$. Equivalently, for every commutative square
\[
\begin{tikzcd}
K\arrow[r,"u"]\arrow[d,"a"']&U\arrow[d]\\
L\arrow[r,"v"']&V
\end{tikzcd}
\]
with $K,L$ finite, there is $t:L\to U$ with $ta=u$; the other triangle is not required to commute. Indeed, the finitely many elements and structure equations of $L$, together with $t(a(k))=u(k)$, give exactly such a system. Conversely, a finite system and its parameters can be encoded by a morphism between finite spans. Write $\mathcal P$ for the class of pure monomorphisms.

\textbf{Proposition (comparison with the manuscript).}\ours
\emph{The class $\mathcal P$ is L-$\pi$-accessible for every infinite cardinal $\pi$ in this filtration. It contains the isomorphisms and is stable under pushouts, composition, transfinite composition, and retracts. Nevertheless there is no set $S\subset\mathcal P$ with $\mathcal P\subset\widetilde S$. Consequently the assertion of Theorem~6.1 is false, even for a presheaf category.}

We give the comparison and closure arguments, and use the published non-generation theorem only at the last step. First, filtered colimits preserve purity. Given a finite system with parameters in the source of a filtered colimit arrow and a solution in its target, represent its parameters and solution at a common stage. The finitely many equations hold at a further common stage. Purity there supplies a solution in that stage's source, which gives the required solution in the colimit. The colimit arrow is a monomorphism because filtered colimits of sets preserve monomorphisms. This argument permits arbitrary transition squares; it is not restricted to diagrams of inclusions.

Next fix a pure inclusion $A\subset B$ and a subarrow $U_0\subset V_0$ of size at most $\pi$. Thus $U_0\subset A$, $V_0\subset B$, and $U_0\subset V_0$, but there is no requirement that $U_0=A\cap V_0$. We construct increasing subarrows $U_n\subset V_n$. For every finite equation system with parameters in $U_n$ that is solvable in $V_n$, choose a solution in $A$, using the purity of $A\subset B$. There are at most $\pi$ systems to consider. Adjoin all their chosen solutions to $U_n$ and close under the two structure maps of the span, obtaining $U_{n+1}$ of size at most $\pi$. Let $V_{n+1}$ be the subspan of $B$ generated by $V_n\cup U_{n+1}$. Its size is still at most $\pi$. At the countable union,
\[
U=\bigcup_{n<\omega}U_n,\qquad V=\bigcup_{n<\omega}V_n,
\]
the inclusion $U\subset V$ is pure. A finite list of parameters and a finite solution in $V$ both occur at some stage, so its required solution was added at the next stage. Hence the pure subarrows of size at most $\pi$ are cofinal among all subarrows of that size. They are also closed under filtered unions with at most $\pi$ members, by the preceding paragraph. Conversely, an arrow whose bounded pure subarrows are cofinal is their filtered colimit and is pure. This is precisely the cofinal small-subobject recognition condition of Proposition~5.7(c), applied to $\underline{\Fl}(\mathcal E)$. That proposition therefore gives the manuscript's L-$\pi$-accessibility, with its original, stronger meaning.

For completeness, pushout stability can be proved without a further smallness assumption. Every pure arrow is a filtered colimit of split monomorphisms between finite spans. To see this, the category of finite arrows mapping to a fixed arrow $f:A\to B$ is filtered, and its colimit is $f$: finite sets of elements and relations can be put into one finite subarrow. If $a:K\to L$ is such an arrow, purity gives $t:L\to A$ extending its map $K\to A$. Form
\[
D=L\amalg_K L,
\qquad j_0,j_1:L\longrightarrow D.
\]
The arrow $j_0$ is split by the fold $D\to L$. There is a morphism $a\to j_0$ in the arrow category with source component $a$ and target component $j_1$. There is also a morphism $j_0\to f$: its source component is $t$, and its target component is induced by $ft$ on the first copy of $L$ and the original $L\to B$ on the second. These two maps agree on $K$. Thus every finite arrow over $f$ maps to a finite split monomorphism over $f$. The full subcategory of such split monomorphisms is filtered: any finite diagram first has a cone in the category of finite arrows over $f$, and the construction just given replaces that cone by a split one. For any fixed finite arrow $a\to f$, the same construction, keeping the map from $a$ as part of the diagram, makes its comma category of split targets nonempty and filtered. This proves finality, so the colimit of the split subcategory is $f$.

Now push out this filtered presentation of $f$ along any $A\to A'$. The resulting arrow is a filtered colimit of pushouts of split monomorphisms, each of which is split. It is pure by filtered-colimit stability. Composition of pure inclusions follows by reflecting a finite system twice. Retract stability follows by composing a finite square with the section into the larger arrow, taking its pure extension, and composing with the retraction. Transfinite composition follows by induction, using composition at successor stages and filtered-colimit stability at limit stages. These arguments establish every closure property asserted above.

The external ingredient is now exact and short. \citet[Theorem~5.1]{CoxFeigertKamsmaMazariArmidaRosicky2026} prove that pure monomorphisms in $\mathrm{Set}^{D}$ are cofibrantly generated if and only if every coslice $d/D$ is linearly preordered. For our $P$, the two objects $0\to1$ and $0\to2$ of $0/P$ admit no morphism in either direction. Thus $\mathcal P$ is not cofibrantly generated. If Theorem~6.1 held, it would give a set $S\subset\mathcal P$ with $\mathcal P\subset\widetilde S$. The small object argument in $\mathcal E$ identifies $\widetilde S$ with the closure of $S$ under pushouts, transfinite composition, and retracts. The closure properties just proved give the reverse inclusion $\widetilde S\subset\mathcal P$. This would make $\mathcal P$ cofibrantly generated, a contradiction.\hfill$\square$

This also explains why the successful theorem of the next section has genuinely more structure. Purity has cancellation in the direction ``$gf$ pure implies $f$ pure''. Grothendieck's construction in Section~7 needs the other cancellation: from $f$ and $gf$ in $W$, it must recover $g\in W$. That step keeps the residual inclusion in the desired class after a small attachment. It is precisely what the preceding argument cannot obtain for a general L-accessible class.

The non-generation theorem and its presheaf classification belong to Cox, Feigert, Kamsma, Mazari-Armida, and Rosick\'{y}; no priority claim is made here. Our contribution is the explicit check of the manuscript's bounded recognition condition and the resulting application to Theorem~6.1. We checked Theorem~5.1 and its reduction to the preceding results, Definition~2.7, Remark~2.8, Fact~2.9, and Example~5.3 in the 26-page third arXiv version, dated 29 April 2026. We did not find this specific application to \emph{Les D\'erivateurs} stated there.
\end{dcomment}

% Sections 7--8. Commentary by GPT-6 Astra (OpenAI, ChatGPT Pro mode).
\begin{dcomment}{exposition}{p85}{What the generation argument must preserve}
Grothendieck now changes the question in a decisive way. The unsuccessful argument of \S6 tried to recover a lifting class from accessibility alone. The useful extra information is that an arrow in the desired class can be enlarged a little while its remaining part stays in the same class. This is a geometric condition on a factorization, and it is stronger than the assertion that the arrow is a filtered colimit of small arrows. A filtered presentation allows both its source and its target to vary. A cellular construction keeps the original source fixed and attaches successive pieces to it. The discrepancy between those two operations was precisely what invalidated the preceding proof.

We use modern lifting notation to make the goal explicit. For a class $S$ of arrows, $S^{\square}$ consists of arrows having the right lifting property against every member of $S$, and ${}^{\square}T$ denotes the corresponding left lifting class. Grothendieck's $S_*$ is $S^{\square}$, and his left $Q$-closure $\widetilde S$ is ${}^{\square}(S^{\square})$. When $S$ is a set and the small object argument applies, this is the closure of $S$ under pushouts, transfinite composition and retracts. A proof that $C\subset\widetilde{C_0}$ must therefore account for this particular closure; being generated by filtered colimits is a different assertion.

The property in Theorem 7.1 supplies the missing control. Suppose $A\hookrightarrow B$ is not yet an isomorphism. A small subobject $B'\hookrightarrow B$ detects something outside $A$, and the union
\[
 D=A\amalg_{A\cap B'}B'
\]
still embeds in $B$. Thus $A\subsetneq D\subset B$ is a genuine enlargement, obtained by one permitted attachment. For trivial cofibrations one must additionally keep $D\hookrightarrow B$ a weak equivalence. This is where cancellation in $W$ enters: once $A\to D$ and $A\to B$ lie in $W$, so does $D\to B$. The hypotheses that look separate in the statement are arranged to secure these two facts, geometric containment and homotopical cancellation.

The question on manuscript page 85 has a precise later answer in the setting of model recognition. If $\mathcal M$ is locally presentable, $C=\operatorname{cof}(I)$ for a set $I$, $W$ is an accessible full subcategory of the arrow category with an accessible inclusion, $W$ has two-out-of-three and is closed under retracts, $I^{\square}\subset W$, and $C\cap W$ is closed under pushouts and transfinite compositions, then $C\cap W$ is generated by a set. This is Smith's theorem, published with proof by \citet[Theorem 1.7 and Lemmas 1.8--1.9]{Beke2000}. Its solution-set hypothesis is supplied by accessibility. It does not assert that the intersection of an arbitrary accessible class with a lifting class is automatically a lifting class: the closure and compatibility hypotheses remain part of the theorem. Consequently it resolves the model-structure problem which motivates this passage, without turning the tentative claim of \S6 into a general theorem.
\end{dcomment}

\begin{dcomment}{examples}{p87}{Small attachments in presheaves and in group actions}
Grothendieck's generator condition has an especially concrete meaning for presheaves. Let $A$ be small and choose an infinite cardinal $\kappa\geq|\operatorname{Mor}A|$. For $F\in\widehat A$, put $|F|=\sum_{a\in\operatorname{Ob}A}|F(a)|$. Up to isomorphism the presheaves of cardinal at most $\kappa$ form a set. Let $M_0$ be a skeleton of this class, and let $I$ contain all monomorphisms between its objects.

Given a proper inclusion $U\subset V$, choose $a$ and $v\in V(a)\setminus U(a)$. The associated map $h_a\to V$ has image $B'$, a subpresheaf of cardinal at most $\kappa$, and $B'$ is not contained in $U$. Its intersection $A'=U\cap B'$ again has cardinal at most $\kappa$. Since unions and intersections are computed objectwise, the canonical map
\[
 U\amalg_{A'}B'\longrightarrow V
\]
is a monomorphism, with image $U\cup B'$. This verifies every part of (7.1.1). The image is essential: $h_a\to V$ need not itself be a monomorphism, so directly attaching $h_a$ need not produce a subobject of $V$. Passing to its quotient image is exactly the reason for allowing quotients among the small pieces.

There is a more explicit calculation when $\widehat A$ is the category of left $G$-sets. Every monomorphism $U\hookrightarrow V$ identifies $U$ with a union of $G$-orbits in $V$, and its complement is a disjoint union of orbits $G/H$. Thus the set
\[
 I_G=\{\varnothing\longrightarrow G/H\mid H\leq G\}
\]
generates all monomorphisms by coproducts and transfinite attachments; retracts are unnecessary for this particular generation. A map $p:X\to Y$ belongs to $I_G^{\square}$ precisely when, for every subgroup $H\leq G$, the map on fixed points
\[
 X^H\longrightarrow Y^H
\]
is surjective. Indeed, an equivariant map $G/H\to Y$ is determined by an $H$-fixed element $y$, and a lift is exactly an $H$-fixed $x$ with $p(x)=y$. Conversely the fixed-point condition solves every orbit attachment and hence every monomorphism lifting problem.

This example explains the word ``injective'' in the source. It refers to an injective object of the slice, tested against all monomorphisms in that slice. Ordinary surjectivity need not suffice. If $G$ is nontrivial, the equivariant map $G\to *$ is surjective as a map of sets, but $G^G$ is empty, so it fails the test for $H=G$. The small object factorization adds precisely the missing equivariant extensions; it does not merely repair the underlying map of sets.
\end{dcomment}

\begin{dcomment}{exposition}{p93}{The transfinite argument and its stopping condition}
The strength of Lemma 7.3 is that its proof never guesses a universal set of cells. It assumes that a permissible next cell can be found whenever a remainder is nontrivial, and then proves that the process reaches the target. Here is the full argument, including its size condition.

Fix $i:A\hookrightarrow B$ in $C$, and let $\mathcal S(B)$ be the set of subobjects of $B$. Choose an ordinal $\theta$ whose cardinal is strictly greater than $|\mathcal S(B)|$. Starting with $A_0=A$, apply condition (b) whenever $A_\alpha\hookrightarrow B$ is not invertible, obtaining a pushout attachment $A_\alpha\to A_{\alpha+1}$ in $\widetilde{C_0}$ and a proper inclusion $A_\alpha\subsetneq A_{\alpha+1}\subset B$ whose remainder lies in $C$. At a nonzero limit ordinal $\lambda$, form $A_\lambda=\varinjlim_{\alpha<\lambda}A_\alpha$. The arrows $A_\alpha\to B$ form a filtered diagram in the arrow category with constant target; condition (c) puts its colimit $A_\lambda\to B$ in $C$, and hence makes it monic. This identifies $A_\lambda$ with the union of the preceding subobjects.

If the construction had not stopped before $\theta$, it would give more distinct subobjects of $B$ than $\mathcal S(B)$ contains. Therefore $A_\lambda=B$ at some stage. Every successor attachment is in $\widetilde{C_0}$, so transfinite closure gives $A\to B\in\widetilde{C_0}$. This also proves the claimed equality of closures when $C_0\subset C$.

The printed phrase asserting the existence of a ``cardinal $\alpha$'' is an ordinal/cardinal slip. The argument produces a stopping \emph{ordinal} $\lambda$ with $|\lambda|\leq|\mathcal S(B)|$, not necessarily a stopping stage which is itself an initial ordinal. The distinction changes no conclusion. More substantively, well-poweredness is explicitly assumed in Lemma 7.3, but is absent from the short list at the start of Theorem 7.1(1). Its proof needs this hypothesis, or another smallness argument that forces termination. It holds in all the presheaf applications here. We retain it when stating the general result; we have not obtained a counterexample to the theorem with that hypothesis omitted.

A second small point explains the retract assertions used subsequently. If $f$ is a retract of $g$ in the arrow category, with section $s$ and retraction $r$, then $e=sr$ is an idempotent endomorphism of $g$. The sequential diagram $g\xrightarrow e g\xrightarrow e\cdots$ has colimit $f$: a cocone is exactly a map out of $g$ invariant under $e$, and these are exactly maps out of its splitting $f$. Thus closure under filtered colimits implies closure under retracts. This observation justifies the source's parenthesis without importing a model structure that has not yet been constructed.
\end{dcomment}

\begin{dcomment}{exposition}{p96}{Keeping the remainder a weak equivalence}
Theorem 7.1(3) succeeds because it repairs the precise defect of \S6. Let $f:A\hookrightarrow B$ lie in $W$. Use the full filtration hypothesis (d) from manuscript pages 45--46, including closure under subobjects; its abbreviated restatement in 6.1 omitted this clause. Write $B$ as the $\pi$-adapted union of its $\pi$-small subobjects $B_i$, and put $A_i=A\cap B_i$. Subobject closure gives $A_i\in\mathrm{Filt}^{\pi}$, and preservation of filtered colimits by pullback gives
\[
 f=\varinjlim_i(A_i\hookrightarrow B_i).
\]
The same preservation proves continuity at the small filtered suprema in the index poset. Thus this is a $\pi$-adapted presentation in the arrow category, rather than merely an arbitrary filtered presentation.

Accessibility now has a legitimate input: it says that the indices for which $f_i:A_i\hookrightarrow B_i$ belongs to $W$ are cofinal. If $A\neq B$, some $B_i$ meets the complement of $A$; cofinality permits choosing such an index with $f_i\in W$. Form
\[
 A\xrightarrow{\alpha_i}D_i=A\amalg_{A_i}B_i
   \xrightarrow{\beta_i}B.
\]
The first arrow is a pushout of a small trivial monomorphism, so it lies both in $\widetilde{(W\cap\mathrm{Mon})^{\pi}}$ and in $W$. Effective unions, condition (f), say that $\beta_i$ is monic. Since $\beta_i\alpha_i=f\in W$, the cancellation clause (7.1.4)(b) gives $\beta_i\in W$. The attachment is proper because $B_i$ was not contained in $A$. These facts verify precisely the successor-step hypothesis of Lemma 7.3 for $C=W\cap\mathrm{Mon}$.

No assertion that the large remainder $D_i\to B$ is itself $\pi$-small has been used. In \S6 that assertion was needed to apply accessibility to a second diagram, and it was unavailable. Here the remainder is put in $W$ by cancellation instead. This is the mathematical reason why the new theorem survives the failure of the earlier one.

Once $TC=W\cap\mathrm{Mon}$ is generated by its small members, the additional closure properties make it equal to its left $Q$-closure. Composition follows from composition in $W$ and in $\mathrm{Mon}$; pushout stability is assumed; transfinite composition and retract closure follow from filtered-colimit stability, with identities supplying the degenerate cases. The small object argument then constructs the weak factorization system $(TC,TC^{\square})$. Its ordinal length must be chosen large enough for the domains of the chosen generating set to be small relative to the construction. Merely saying that the generating set is small is not itself this domain-smallness assertion; the accessibility hypotheses supply it.
\end{dcomment}

\begin{dcomment}{exposition}{p98}{Why the two factorization systems fit together}
The remaining issue in Corollary 7.4 is compatibility. Two weak factorization systems on the same category do not by themselves make a model structure. Here the required compatibility is concentrated in the inclusion $TF\subset W$, and the retract argument explains why this one inclusion is enough.

Assume that $(C,TF)$ and $(TC,F)$ are weak factorization systems, that $TC=C\cap W$, that $TF\subset W$, and that $W$ has two-out-of-three and is closed under retracts. Because $TC\subset C$, every map right orthogonal to $C$ is right orthogonal to $TC$; hence $TF\subset F\cap W$. Conversely, let $p:X\to Y$ lie in $F\cap W$. Factor it using the first system as
\[
 X\xrightarrow{i}Z\xrightarrow{q}Y,
 \qquad i\in C,\quad q\in TF.
\]
Since $q,p\in W$, two-out-of-three gives $i\in W$, hence $i\in TC$. Lifting $i$ against $p$ gives $r:Z\to X$ with $ri=1_X$ and $pr=q$. In the arrow category the maps $(i,1_Y)$ and $(r,1_Y)$ exhibit $p$ as a retract of $q$. Right lifting classes are retract-closed, so $p\in TF$. Consequently
\[
 TF=F\cap W,\qquad TC=C\cap W,
\]
and the two factorizations and lifting properties are exactly the model-category axioms. Completeness and cocompleteness are supplied by the ambient hypotheses in the applications. This proves the sufficiency left implicit in the reference to \S4.3.

For $C=\mathrm{Mon}$ the class $TF$ is independent of $W$. It is already fixed by the ordinary category, and condition (7.2.7) asks the proposed homotopy theory to regard all those maps as equivalences. This is a substantial condition. In $\mathrm{Set}$, for example, $\mathrm{Mon}^{\square}$ is the class of surjections, assuming the usual axiom of choice: testing $\varnothing\to\{*\}$ gives surjectivity, and choices of preimages solve the general lifting problem. If $W$ consists only of bijections, $TF\subset W$ fails. Thus the category of sets with weak equivalences the bijections cannot have all injections as cofibrations and the factorization systems prescribed here.

Left properness has a separate role. It asks that a pushout of any weak equivalence along a cofibration be a weak equivalence. At this stage the source uses its previously established identity $\mathrm{Cof}_W\cap W=W^{\mathrm{univ}}$, where the superscript denotes universal stability under \emph{cobase} change, to obtain pushout stability of $TC$. In the presheaf applications, once the model structure is constructed, there is also the general explanation that every object is cofibrant: initial maps are monomorphisms, and these are the cofibrations. Homotopy invariance of pushouts of cofibrant objects then gives left properness. The assertion about initial maps is a property of presheaves, not of arbitrary categories. None of these arguments gives right properness automatically. The symbols $W^{\mathrm{univ}}$ and $W_{\mathrm{univ}}$ must therefore retain their opposite variances throughout this passage.
\end{dcomment}

\begin{dcomment}{examples}{p99}{Which vertical map is the fibration}
Grothendieck wants the last clause of Corollary 7.4 to record right properness. In the displayed cartesian square the fibration hypothesis belongs to $f':X'\to Y'$, the arrow \emph{along which} $p:Y\to Y'$ is pulled back. The edition instead prints $f\in F$, referring to its pullback. This is a source slip, not a transcription error. The correction is necessary even in the most familiar proper model category.

In simplicial sets with the Kan--Quillen structure, take $Y'=\Delta[1]$, let $Y=\{0\}$, and let $X'=\{1\}$. Both maps to $Y'$ are the indicated endpoint inclusions. Their pullback is empty, so the square is
\[
\begin{tikzcd}
\varnothing\arrow[r]\arrow[d] & \{1\}\arrow[d]\\
\{0\}\arrow[r] & \Delta[1].
\end{tikzcd}
\]
The bottom map is a weak equivalence. The left map $\varnothing\to\{0\}$ is a Kan fibration: every horn is nonempty, so there are no lifting problems with a map from a horn into its empty domain. The top map is not a weak equivalence, since it changes the set of connected components from empty to a singleton. The right endpoint inclusion is not a Kan fibration; a one-dimensional horn asking for the missing endpoint has no lift. Thus the erroneous condition on the left-hand vertical map does not express right properness. Replacing it by $f'\in F$ restores the intended statement.
\end{dcomment}

\begin{dcomment}{development}{p102}{Two finite categorical obstructions to the exactness hypotheses}
Grothendieck separates $\mathrm{Cat}$ from a presheaf topos because the generation argument needs effective unions and stable monomorphisms, not only filtered-colimit exactness. His doubts about these two properties in $\mathrm{Cat}$ admit very small witnesses. We give them explicitly, so that the restriction to presheaves can be seen in the actual constructions. We checked the model-recognition discussions of \citet[Section 1]{Beke2000} and \citet[Sections 1.2 and 4.2]{Cisinski2006}; neither is being cited as the source of the particular finite presentations below. The examples establish familiar failures and carry no claim of priority.

\textbf{Failure of effective unions.}\ours\ Let $B$ be the four-element diamond poset, with objects $0,1,2,3$ and nonidentity relations
\[
 0<1<3,\qquad 0<2<3.
\]
Let $A_1$ contain all four objects and just the nonidentity arrows $0\to1$ and $2\to3$, and let $A_2$ contain all four objects and just $0\to2$ and $1\to3$. Each is a subcategory of $B$, and $A_0=A_1\cap A_2$ is discrete. Their pushout $P=A_1\amalg_{A_0}A_2$ is the free category on the directed diamond: a functor from $P$ is exactly the data of functors from $A_1$ and $A_2$ agreeing on objects, with no imposed relation between the two length-two paths. Thus $P$ has two distinct arrows $0\to3$, whereas $B$ has one. The comparison $P\to B$ is not faithful and therefore is not a monomorphism in $\mathrm{Cat}$. This disproves condition (f) of (7.2.1).

\textbf{Failure of stability of monomorphisms under pushout.}\ours\ Let $A=[1]$, with arrow $a:0\to1$, and let $B$ be the category obtained by making $a$ invertible. The functor $A\to B$ is injective on objects and faithful, hence monic. Let $C$ have objects $0,1,2$, arrows
\[
 a:0\to1,\quad b,c:1\to2,\quad d:0\to2,
 \qquad ba=d=ca,
\]
with $b\neq c$ and no further nonidentity arrows or relations. This is a finite category, and $A\to C$ is the evident inclusion. In the pushout $D=B\amalg_A C$, the image of $a$ is invertible. Therefore the equality $ba=ca$ forces $b=c$ in $D$. The induced functor $C\to D$ is not faithful. It is the pushout of the monomorphism $A\to B$, so condition (g) also fails.

Both failures concern creation and identification of composites. In a presheaf category the corresponding constructions are computed in sets at each object; no new paths have to be composed, and there is no cancellation forced by an inverse newly adjoined in another piece. For the finite-union example the precise defect is an additional relation on two newly composable paths. For the pushout example it is a cancellation relation imposed by a new inverse. These are different obstructions, which is why conditions (f) and (g) were listed separately.
\end{dcomment}

\begin{dcomment}{exposition}{p103}{The filtration and the inverse-image lemma}
Grothendieck needs the class $W_A$ to inherit the smallness properties of $W$ before the generation theorem can be applied. The cardinal bookkeeping is straightforward once the correct presentation is fixed. For an infinite $\pi$, a category with at most $\pi$ arrows has at most $\pi$ objects. A colimit of a diagram with at most $\pi$ arrows, whose values have at most $\pi$ arrows, is a quotient of the free category on a graph of size at most $\pi$. Finite strings in such a graph still form a set of size at most $\pi$, so the same cardinal bound holds for the colimit. Finite limits are subsets of finite products of object and arrow sets, with composition inherited componentwise. The estimate for larger limits when $\pi=2^c$ uses $(2^c)^c=2^c$, for infinite $c$.

Every set of at most $\pi$ arrows of an arbitrary category is contained in a subcategory with at most $\pi$ arrows, obtained by adding identities and finite composites. The poset of these small subcategories is $\pi$-directed in the source's cardinal convention, admits the required small directed unions, and has colimit the whole category. If the ambient arrow set has cardinal $\pi'$, there are at most $(\pi')^{\pi}$ such subcategories. For presheaves, restriction maps replace finite composites in the same calculation, and the bound $\pi\geq|\operatorname{Mor}A|$ controls the closure under all restrictions.

Lemma 8.1 then amounts to transporting an adapted presentation. Choose $\pi$ for which $M'$ is $\pi$-accessible and $\varphi(\mathcal M^{\pi})\subset\mathcal M'^{\pi}$. For a $\pi$-adapted presentation $X=\varinjlim_iX_i$, preservation of filtered colimits gives an adapted presentation $\varphi X=\varinjlim_i\varphi X_i$. Membership of $\varphi X$ in $M'$ is equivalent to cofinality of the indices with $\varphi X_i\in M'$. This is exactly the criterion for $X\in\varphi^{-1}M'$. Closure of the inverse image under arbitrary filtered colimits follows separately from the same preservation property. The proof establishes the source's $L$-accessibility; it does not identify this notion merely by terminology with modern accessibility of a full arrow subcategory.

For $i_A:F\mapsto A/F$, the objects are pairs $(a,x)$ and an arrow $(a,x)\to(b,y)$ is an arrow $u:a\to b$ with $x=F(u)(y)$. Filtered colimits of these finite data are computed as expected, so $i_A$ preserves them. Its cardinal bound follows from the number of choices of $u$ and $y$. On enlarging the initial cardinal, the inverse-image lemma therefore applies to $W_A=i_A^{-1}W$.
\end{dcomment}

\begin{dcomment}{status}{p104}{The forward gluing theorem and the extra converses}
Grothendieck's immediate objective is a model structure, for which the forward implications in (8.2.4) are sufficient: pushouts of weak equivalences along monomorphisms remain weak equivalences, and pushouts of trivial monomorphisms remain trivial. The two converse implications printed in that formula impose a stronger requirement. They belong to the strengthened excision framework invoked here and must not be silently identified with the ordinary gluing theorem for homotopy types.

This distinction has already been analysed in the commentary on Chapter~XII, manuscript p.~78 (``A finite counterexample to the asserted reverse excision'') and p.~126 (``Why converse excision is an additional axiom''). We refer to those blocks rather than reproduce their examples. In particular, an application to $W_\infty$ must retain the forward statements and use the actual presheaf gluing theorem, rather than assume that the stronger biconditionals hold for ordinary weak homotopy equivalences.

The surviving construction is substantial. Under the forward stability, the two weak factorization systems in Corollary 8.3 exist, and the sole remaining model-recognition condition is that maps right orthogonal to all monomorphisms belong to $W_A$. Theorem 8.4 turns that apparently large condition into an interval condition. Later, \citet[Corollary 4.2.18]{Cisinski2006} proves the model-structure theorem for every accessible fundamental localizer and every corresponding local test category. This theorem supplies the intended application without the reverse-excision assertion.
\end{dcomment}

\begin{dcomment}{exposition}{p106}{The Lawvere object and the local interval test}
Grothendieck is looking for an intrinsic cylinder in $\widehat A$. The subobject classifier $\Omega=L_A$ is available before any homotopy theory is chosen, and its two sections classify the empty and the full subobject. They are disjoint because a sieve cannot be both empty and maximal. Thus $\Omega$ is always a separated interval as a piece of topos theory; the additional homotopical question is whether its projection to the terminal object is universally a weak equivalence.

Injectivity of $\Omega$ has a direct proof. Given a monomorphism $U\hookrightarrow V$ and a map $U\to\Omega$, let $S\hookrightarrow U$ be its classified subobject. The composite $S\hookrightarrow V$ is a subobject of $V$, so its characteristic map $V\to\Omega$ extends the original map: pulling $S\hookrightarrow V$ back to $U$ recovers $S\hookrightarrow U$. Consequently $\Omega\to1$ has the right lifting property against all monomorphisms. The same is true after every base change, since right lifting properties are stable under pullback. This proves the implication (a)$\Rightarrow$(b) in 8.4, including its universal quantifier.

Checking universal asphericity on representables is the relevant local criterion. For $a\in A$, the slice topos $\widehat A/h_a$ identifies with $\widehat{A/a}$. Under that identification $\Omega\times h_a\to h_a$ is its subobject classifier. More generally, a separated interval $I$ with $I\times h_a\to h_a$ a weak equivalence for every $a$ gives a cylinder after passage to every slice. The localizer's localization axiom then extends the test to arbitrary bases. The equality $L_A\times a=L_{A/a}$ in the source is understood in this slice equivalence; the two objects are not literally presheaves on the same category before that identification.

The category-of-elements functor also explains the right adjoint appearing in (8.4.2). A map $F\to j_A(C)$ assigns to every element $x\in F(a)$ a functor $A/a\to C$, compatibly with restriction along arrows of $A$. This is equivalent to a functor $A/F\to C$: evaluate the assigned functor at $1_a$, and use the compatibility to recover its action on arrows. Conversely, restrict a functor $A/F\to C$ along the canonical functor $A/a\to A/F$ attached to $x$. These constructions are inverse. Hence $i_A\dashv j_A$ is an ordinary adjunction, and the stronger test-category assertion says that its unit and counit become weak equivalences on all relevant slices.

The interval characterization is already present in \citet[\S65, Proposition 3 and Corollary 1]{Grothendieck1983}; in the 2015 transcription the proposition spans printed pages 171--172, corresponding to the manuscript pagination cited here. It is the local condition, followed by the additional asphericity of $A$, which distinguishes a local test category from a test category.
\end{dcomment}

\begin{dcomment}{examples}{p107}{A terminal category fails and a separated interval succeeds}
Theorem 8.4 is not saying that every category with a terminal object is a test category. For $A=e$, presheaves are sets and $i_A$ regards a set as a discrete category. For $W=W_\infty$, the induced class $W_A$ is therefore the bijections. The Lawvere object is the two-element set, whose projection to a singleton is not a weak equivalence. Equivalently, there cannot be a separated interval in sets which is weakly equivalent to a point. The terminal object of $A$ supplies asphericity of the indexing category, but cannot supply the missing interval.

The positive criterion in Corollary 8.5(b) fills exactly this gap. Suppose $A$ has a terminal object and an object $I$ with disjoint endpoint sections, and that $a\times I$ exists for every $a$. The representable presheaf $h_I$ has disjoint sections, and
\[
 h_a\times h_I\simeq h_{a\times I}.
\]
Its category of elements is $A/(a\times I)$, which has a terminal object. Thus the required local asphericity follows from the terminal-object axiom. Since $A$ itself has a terminal object, the additional global asphericity holds as well. This proves the weakened finite-product hypothesis in footnote 40; no products unrelated to $I$ were used.

For the ordinary simplex category, the interval is $\Delta[1]$. The product $\Delta[n]\times\Delta[1]$ is the nerve of the product poset $[n]\times[1]$, whose largest element contracts its nerve. The two endpoint simplicial subsets are disjoint. This supplies the concrete homotopy behind the local criterion, while the terminal object $[0]$ makes $\Delta$ aspherical. The names of the cubical categories in the source require their prescribed choices of maps; one should not transfer the statement to every category called ``cubical'', since diagonals, symmetries and connections change the indexing category.

Footnote 38 distinguishes a genuinely stronger, false reformulation. In the simplicial case let $u:e\to[1]$ select one endpoint. Both categories have terminal objects, so $u\in W_\infty$. The two endpoint maps $j_A(e)\to j_A([1])$ are disjoint: at every $a$ they are the distinct constant functors $A/a\to[1]$. Pulling one endpoint back along the other gives $\varnothing\to j_A(e)=1$. For $A=\Delta$ this is not a weak equivalence. Hence $j_A(u)$ is not universally a weak equivalence, although $A$ is a test category. The weaker condition involving only projections $C\to e$ with $C$ terminal-pointed avoids this counterexample, exactly as the editors' bracketed correction says.
\end{dcomment}

\begin{dcomment}{exposition}{p109}{The retract in Lemma 8.6}
Grothendieck isolates a useful recognition principle: to prove that a right lifting class consists of weak equivalences, it can be enough to factor its arrows through weak equivalences. Write $f:X\to Y$ in $TF$ as $f=pi$, where $i:X\to Z$ is in $C$ and $p:Z\to Y$ is in $W$. The lifting square with top map $1_X$, left edge $i$, bottom map $p$, and right edge $f$ yields $r:Z\to X$ satisfying
\[
 ri=1_X,\qquad fr=p.
\]
Therefore $(i,1_Y):f\to p$ and $(r,1_Y):p\to f$ exhibit $f$ as a retract of $p$ in the arrow category. Retract closure of $W$ gives $f\in W$. The final $i\in W$ printed in the proof is a slip: the assertion proved and needed is $f\in W$.

The other direction uses the factorization already included in the precouple hypothesis: factor any arrow as a $C$-map followed by a $TF$-map, and then use $TF\subset W$. Thus one must retain both ingredients of the precouple, factorization and the lifting relation; a bare orthogonality inclusion would not prove (a)$\Rightarrow$(b). No two-out-of-three property is needed for this lemma. Its economy is what makes it attractive here, where the interval construction is intended to provide the weaker factorization $C$ followed by $W$ before the model structure is known.
\end{dcomment}

\begin{dcomment}{exposition}{p112}{What the mapping cylinder proves}
The mapping cylinder converts a homotopy between two copies of the source into a factorization of an arbitrary map. For $f:X\to Y$ in a presheaf category, form
\[
 Z(f)=Y\amalg_{X,\delta_0}(X\times I),
 \qquad i_f(x)=[(x,\delta_1)],
\]
and let $p_f:Z(f)\to Y$ equal $1_Y$ on $Y$ and $f\,\mathrm{pr}_X$ on $X\times I$. The compatibility on the identified copy of $X$ makes $p_f$ well defined, and $p_fi_f=f$.

There are two independent verifications. First, the map
\[
 Y\amalg X\longrightarrow Z(f)
\]
is the pushout of the boundary inclusion $X\amalg X\hookrightarrow X\times I$. It is monic, as are coproduct inclusions in a presheaf category. Their composite $i_f:X\to Y\amalg X\to Z(f)$ is therefore monic. This proves the assertion about the left factor and explains why disjoint endpoints are needed.

Second, $\mathrm{pr}_X\in W_A$, and its sections $\delta_0^X,\delta_1^X$ lie in $W_A$ by two-out-of-three. If trivial monomorphisms are stable under pushout, then the map $\beta:Y\to Z(f)$ is in $W_A$. Since $p_f\beta=1_Y$, two-out-of-three gives $p_f\in W_A$. This proves the assertion about the right factor. The two arguments use different hypotheses: ordinary monomorphism stability for the first, and \emph{weak-equivalence} stability for the second.

The construction is functorial in maps $Y\to Y'$ under the fixed source $X$, because the universal property of the pushout supplies $Z(f)\to Z(f')$ and these maps respect identities and composition. A varying-source functorial factorization additionally requires a functorial choice of the cylinders $I_X$ and their structure maps. In the present application $I_X=X\times I$, so this further functoriality is available. Compare the commentary on Chapter~XII, manuscript p.~230 (``The mapping cylinder, with the two endpoints distinguished''); the present discussion addresses the additional issue of which hypotheses remain when the argument is extracted as Lemma 8.7.
\end{dcomment}

\begin{dcomment}{development}{p113}{The omitted hypothesis in the abstract cylinder lemma}
Grothendieck has established the mapping-cylinder argument in a presheaf category where trivial monomorphisms are stable under pushout. Lemma 8.7 then extracts an abstract statement but retains only stability of $C$ under composition and pushout. That suffices to prove $i_f\in C$ under the stated coproduct condition; it does not suffice to prove $p_f\in W$. The following example separates the two assertions completely. We checked the cylinder criterion of \citet[Lemma 1.4.13]{Cisinski2006} and the model-recognition argument of \citet[Section 1]{Beke2000}; we did not find this particular counterexample in those sources.

\textbf{Counterexample to Lemma 8.7 as stated.}\ours\ Take $\mathcal M=\mathrm{Grp}$, let $C$ be all group homomorphisms, and let $W$ consist of the homomorphisms $u:G\to H$ such that
\[
 \operatorname{Hom}(C_2,G)\longrightarrow
 \operatorname{Hom}(C_2,H)
\]
is a bijection. The class $W$ is the inverse image of the isomorphisms of sets under a functor. It therefore contains all isomorphisms, has two-out-of-three, and is closed under retracts; it also satisfies the usual saturation condition involving two opposite composites. The category has finite colimits, and $C$ satisfies every stability and coproduct condition in the lemma.

Let $X=1$, the trivial group, and choose $I_X=\mathbb Z$. There is a factorization of the fold map
\[
 X*X=1\longrightarrow\mathbb Z\longrightarrow1=X.
\]
Its left factor is in $C$, and its right factor is in $W$, since $\mathbb Z$ has no nontrivial element of order two. Both endpoint maps are the unique map $1\to\mathbb Z$; the abstract lemma does not require them to be disjoint, and its condition on $i_X$ is satisfied.

For $f:1\to C_2$, the prescribed mapping cylinder is the free product
\[
 Z(f)=C_2*\mathbb Z,
 \qquad p_f:C_2*\mathbb Z\longrightarrow C_2,
\]
where $p_f$ is the identity on $C_2$ and kills the generator $z$ of $\mathbb Z$. If $t$ is the nonidentity element of $C_2$, both $t$ and $ztz^{-1}$ are involutions in the free product. They are distinct by reduced-word normal form, but $p_f$ maps both to $t$. Thus the induced map on $\operatorname{Hom}(C_2,-)$ is not injective, and $p_f\notin W$. Every explicit hypothesis of Lemma 8.7 holds, while its right-factor conclusion fails.

The repair can be stated at the exact level used by the proof. Assume that $\delta_0:X\to I_X$ belongs to $W$ and that its pushout along each map $f:X\to Y$ under consideration belongs to $W$. Then $\beta:Y\to Z(f)$ is in $W$, and $p_f\beta=1_Y$ gives $p_f\in W$. A uniform sufficient condition is $C\cap W$ stable under pushout, together with $\delta_0\in C$. The latter follows if $X\to X\amalg X$ and $i_X:X\amalg X\to I_X$ are in $C$. For the left factor one also needs $X\to Y\amalg X$ in $C$, precisely the explicit condition used in the printed proof. In the usual convention this is automatic when $Y$ is cofibrant, and it holds for every $X,Y$ in a presheaf category with $C=\mathrm{Mon}$. Calling it simply ``$X$ cofibrant'' reverses which added summand supplies the condition.

With these hypotheses the proof on the preceding page is complete: $j:Y\amalg X\to Z(f)$ is a pushout of $i_X$, so $i_f=j\circ(X\to Y\amalg X)$ lies in $C$; the separate argument with $\beta$ proves $p_f\in W$. The failure of the abstract lemma therefore does not invalidate the presheaf application under its original gluing assumptions.
\end{dcomment}

\begin{dcomment}{exposition}{p114}{The interval proof that avoids pushouts}
Grothendieck's final reduction of the hypotheses in 8.4 asks for an argument that no longer uses the pushout axiom. Such an argument exists: the map under examination already has a lifting property, so the interval can be used to construct a homotopy directly. This is the proof later given in \citet[Lemma 1.4.13, printed pp.~56--57]{Cisinski2006}. We spell it out with the hypotheses of the present text.

Suppose $\widehat A$ has a separated interval $I$ such that $X\times I\to X$ belongs to $W_A$ for every $X$, and suppose $W_A$ has two-out-of-three and is closed under retracts. Let $p:X\to Y$ have the right lifting property against all monomorphisms. Lifting against $\varnothing\hookrightarrow Y$ gives a section $s:Y\to X$, with $ps=1_Y$. The endpoint inclusion
\[
 X\amalg X\hookrightarrow X\times I
\]
is monic. Map its two summands to $X$ by $1_X$ and $sp$, and map $X\times I$ to $Y$ by $p\,\mathrm{pr}_X$. The equalities $p1_X=p$ and $psp=p$ make this a commutative lifting square. A lift $H:X\times I\to X$ therefore satisfies
\[
 H\delta_0=1_X,\qquad H\delta_1=sp.
\]
Both endpoint maps are in $W_A$, because they are sections of $\mathrm{pr}_X\in W_A$. From $H\delta_0=1_X$, two-out-of-three gives $H\in W_A$, and hence $sp=H\delta_1\in W_A$.

Finally $p$ is a retract of $sp$ as an arrow: the maps are $(1_X,s):p\to sp$ and $(1_X,p):sp\to p$, and their composite on $p$ is $(1_X,ps)=1_p$. Retract closure gives $p\in W_A$. This proves (c)$\Rightarrow$(a) without a single pushout, without accessibility, and without a model structure. If one uses the weak-saturation axiom saying that $ps$ and $sp$ in $W_A$ force $p,s\in W_A$, the final retract step may be replaced by that axiom.

The reverse implications remain the elementary ones: the subobject classifier is injective, and its empty/full sections form a separated interval. Thus the equivalence of (a), (b) and (c) really has the weaker character claimed at the end of \S8. The argument corrects a proof dependency, not the theorem's conclusion. Accessibility is needed later to produce factorizations; it is not needed to decide whether this interval test recognizes the intended weak equivalences.
\end{dcomment}

\begin{dcomment}{later}{p115}{What later model recognition establishes}
The construction sought here became a theorem with fewer auxiliary hypotheses. For a fundamental localizer $W$ generated by a set and a $W$-local test category $A$, \citet[Corollary 4.2.18]{Cisinski2006} constructs the cofibrantly generated model structure with weak equivalences $W_A$ and cofibrations all monomorphisms. The interval characterization is Theorem 4.1.19. The source's question about right properness is a separate question: Cisinski's Theorem 4.3.24 makes properness of the fundamental localizer independent of the chosen test category, and the minimal localizer is proper (Example 4.3.22).

There are three smallness statements to keep distinct. Grothendieck's $L$-accessibility includes closure under \emph{every} filtered colimit together with his adapted-presentation criterion. Modern accessibility of a full subcategory of the arrow category concerns a sufficiently large regular cardinal and an accessible inclusion. Cisinski calls a localizer accessible when it is generated by a \emph{set as a localizer}; for presheaf localizers, his Scholie 1.4.6 compares this with accessibility of the arrow class. These comparisons are theorems, not interchangeable definitions. In particular, the model-recognition theorem for a general accessible presheaf localizer should not be reported as assuming closure under all filtered colimits merely because the present chapter does so.

The logical order of the achievement also matters. A separated interval first proves that injective-type maps are weak equivalences; a bounded generation argument then supplies the weak factorization systems; their compatibility yields the model structure. The associated derivator can subsequently be formed from diagram categories and their weak equivalences. This order explains why the elementary topos-theoretic object $L_A$ is central to a chapter whose ultimate purpose is the construction of homotopy theories.
\end{dcomment}

% Section 9--10 commentary, GPT-6 Astra Pro (OpenAI).
\begin{dcomment}{exposition}{p116}{The catalogue distinguishes several notions of localizer}
Grothendieck is reorganizing the hypotheses so that each construction can carry only the assumptions it needs. The list is a catalogue of possible conditions, not a declaration that all of them are in force. This matters particularly for localization: in these pages a ``localiseur fondamental'' without qualification satisfies Loc\,(1,2,3), whereas the subsequent standard definition uses the relative condition Loc\,(3\,bis).

The distinction can be expressed without terminology. The weak condition recognizes a functor $u:A\longrightarrow B$ when every comma category $A/b$ has the homotopy type of a point, relative to $W$. The relative condition starts with a map of diagrams over a third category $S$ and recognizes $u$ from the maps $A/s\longrightarrow B/s$. Here $A/s$ has objects $(a,p(a)\longrightarrow s)$; it is not generally the strict fibre of $p:A\longrightarrow S$. The equivalence between the two formulations of Loc\,(3) uses the terminal object $(b,1_b)$ of $B/b$ and two-out-of-three. It does not make the relative condition redundant.

There is a similarly necessary distinction within ``saturation''. Loc\,(1\,c) is the two-sided inverse test: if $gf$ and $fg$ belong to $W$, then so do $f$ and $g$. Loc\,(1\,bis) adds closure under retracts in the arrow category. Strong saturation, later called Loc\,(1\,ter), says that an arrow belongs to $W$ precisely when its image in $\Cat[W^{-1}]$ is invertible. The last condition implies the preceding ones, but its proof requires work; it is not part of the three displayed clauses of Loc\,(1).

The modern vocabulary records these distinctions. The definitions of weak and ordinary fundamental localizers in \citet[\S\S1.1 and 2.1]{Maltsiniotis2020} use respectively the absolute and relative forms of Theorem A. The theorem that every weak fundamental localizer contains the nerve weak equivalences is \citet[Theorem 6.1.18]{Cisinski2006}. Containment of the usual weak equivalences does not, by itself, supply every stability assertion concerning the additional arrows of $W$.

We should also keep the strength of Loc\,(4) visible. Its reverse implication is an extra excision condition, not a general property of homotopy pushouts. The commentary on Chapter~XII already gives the relevant warnings: manuscript p.~126 (``Why converse excision is an additional axiom'') and, in the appendix, manuscript p.~17 (``Acyclic links and a false converse''). Those examples are not repeated here. In this chapter Loc\,(4) may legitimately be assumed; it should not be silently assigned to every localizer obtained from a familiar homotopy theory. Finally, the parenthetical reduction of Loc\,(3\,bis) to split fibred categories is proposed rather than demonstrated. A proof must specify its fibration convention and compare the comma resolutions with the original categories; replacing comma categories by strict fibres without that comparison would lose the hypothesis being tested.
\end{dcomment}

\begin{dcomment}{context}{p117}{Which fibration hypothesis is being invoked}
The variants imported from chapter XII concern the old notion of $W$-fibration. Write $F_{\mathrm{old}}$ for that class, reserving $\Fib_W$ for the intrinsic condition that pullback preserves every $W$-map. The reference is \citet[ms.~53]{Grothendieck1990XIIa}; after the renumbering, its W\,(4) becomes Loc\,(6). The seven tests are as follows.
\[
\begin{array}{c|l}
 a&F_{\mathrm{old}}\subseteq\Fib_W;\\
 a'&\text{the same test restricted to }\mathrm{Parf}_{\omega};\\
 b&f:X\to Y\text{ over }S,\ p,q\in F_{\mathrm{old}},\ f\in W
       \ \Longrightarrow\ f_s\in W\text{ for all }s;\\
 b'&\text{the preceding test for split fibred objects and split cartesian }f;\\
 b''&\text{the preceding test with }p,q\text{ both }W\text{-perfect};\\
 c&p\in F_{\mathrm{old}}\cap W\ \Longrightarrow\ p^{-1}(s)\text{ is }W\text{-aspherical};\\
 c'&\text{the preceding test restricted to }\mathrm{Parf}_{\omega}\cap W.
\end{array}
\]
Here $p:X\to S$, $q:Y\to S$, and $\mathrm{Parf}_{\omega}$ refers to the least weak fundamental localizer. The restricted classes are part of the hypotheses. In particular, the proof on ms.~121 invokes $c'$ after writing only $p\in\mathrm{Parf}_W$. It needs a replacement in the restricted perfect class, together with componentwise weak equivalences for its first factor. These are dependencies of that short argument. The modern theorem identifying the least weak fundamental localizer with $W_\infty$ clarifies the subscript, but does not supply an unmentioned replacement construction merely by changing the notation.
\end{dcomment}

\begin{dcomment}{context}{p118}{A size condition with a particular closure operation}
The accessibility hypothesis is doing a different job from the homotopical axioms. It is intended to turn a proper class of potential attachments into a set from which a factorization can be constructed. In Loc\,(7\,bis), the closure operation is specifically closure under all small filtered colimits in the arrow category of $\Cat$. Its relationship with the cardinal filtration was developed in \S5. Thus ``L-accessible'' contains both a bounded generation assertion and a closure assertion.

Later authors also call a fundamental localizer accessible when a set generates it \emph{as a fundamental localizer}; this is \citet[Definition 4.2.10]{Cisinski2006}. The latter closure includes the localizer axioms. These two formulations must not be identified simply because they share the word ``accessible''. A comparison requires a theorem connecting the two closure operations. The theorem on filtered colimits in the mature theory is recorded separately in \citet[Proposition 2.4.12 and Corollary 2.4.14]{Maltsiniotis2020}. The size and closure requirements in Grothendieck's stated theorems are retained in the audit rather than replaced by an unqualified modern accessibility assumption.
\end{dcomment}

\begin{dcomment}{exposition}{p120}{Adding a summand and recovering a summand}
The point of this discussion is to make weak equivalence compatible with the decomposition of a category into connected components. There are two operations to control. One adds an unrelated component to an equivalence. The other extracts a component from an equivalence already given. Proposition 9.4 separates them cleanly, and its proof is entirely formal once those operations are stated.

Write $P$ for the condition that $f\in W$ implies $f\amalg1_Z\in W$, and $R$ for the condition that restriction to a summand of the target preserves $W$. Assuming identities and composition belong to the saturation rules, $P$ implies closure under binary sums because
\[
 f_1\amalg f_2=(1_{Y_1}\amalg f_2)(f_1\amalg1_{X_2}).
\]
Conversely, binary closure implies $P$. If $f_1\amalg f_2\in W$, condition $R$, applied to each of the two target summands, recovers $f_1$ and $f_2$. If binary sums are both preserved and reflected, any $f:X\longrightarrow Y_0\amalg Y_1$ decomposes canonically as $f_0\amalg f_1$, since a functor cannot carry an arrow between two different components of the target. Reflection then proves $R$.

The same argument works for an arrow $X\longrightarrow X'$ over a category $Y$ and a summand $Y_0\subset Y$: decompose $X$ and $X'$ by inverse image of $Y_0$ and its complement. This proves that $R$ is exactly the assertion that the inclusion $Y_0\hookrightarrow Y$ belongs to $\Fib_W$. The target of the restricted arrow in Corollary 9.3 must therefore be $X'_0$, not $Y_0$. The latter is the base of the restricted diagram, not its target.

Reflection for an arbitrary family already follows from binary reflection. For a fixed index $i$, group all other indices into one summand. No limiting argument is needed for this direction. Preservation for an infinite family requires an additional step: the coproduct arrow is the filtered colimit of its finite subcoproducts, indexed by the finite subsets of the index set. Thus Loc\,(7) upgrades finite preservation to arbitrary preservation. This explains why Loc\,(5\,bis) only adds the infinite version of the forward implication. It also proves the two equivalent componentwise formulations of Corollary 9.5.

The proof of Corollary 9.2 deserves a small qualification. Its displayed square is a pushout along a summand inclusion, but the horizontal arrow $f$ is arbitrary. Literally, Loc\,(4) on ms.~117 is phrased for a union of two open subcategories. To apply an excision theorem to this square one must use the corresponding extension to pushouts along Dwyer inclusions, or first perform the mapping-cylinder reduction established in the preceding discussion. The word ``immédiat'' suppresses that extension. Under the stronger pushout formulation the claimed equivalence is immediate; the displayed square alone is not an instance of the literal two-open-cover formulation. We do not infer a counterexample to the conclusion from this mismatch of formulations.
\end{dcomment}

\begin{dcomment}{development}{p125}{Connectivity supplies the missing retraction}
Grothendieck asks whether the decomposition axiom and the component axiom are independent. The question changes once closure under retracts is available. The following elementary observation makes the relation precise, and accounts for why the modern situation is more rigid than the tentative list suggests. We checked the localizer discussions in \citet[\S2.1 and \S2.4]{Maltsiniotis2020} and \citet[\S9.3]{Cisinski2006}; the particular retract argument below is supplied here, without a claim of priority.

\textbf{Summand reflection from connectivity.}\ours
Let $W$ be a class of functors containing the identities and closed under retracts in $\uFl(\Cat)$. Suppose every $W$-map induces a bijection on connected components. Then, for every small family $(f_i:X_i\longrightarrow Y_i)_{i\in I}$,
\[
 \coprod_{i\in I}f_i\in W\quad\Longrightarrow\quad f_i\in W
 \quad\text{for every }i.
\]
Consequently, if $W$ is also closed under arbitrary coproducts, it satisfies Loc\,(5\,bis).

\emph{Proof.} Fix $i$. If $X_i$ is empty, then $Y_i$ is empty: otherwise the component of any object of $Y_i$ would not be in the image of $\pi_0(\coprod f_j)$, since all the other source summands map into the other target summands. In this case $f_i$ is an identity. If $X_i$ is nonempty, choose an object $x_i$. The inclusion of $f_i$ into $\coprod f_j$ is an arrow-category map. Define a retraction on sources to be the identity on $X_i$ and, on each $X_j$ with $j\ne i$, the constant functor at $x_i$. Define the retraction on targets to be the identity on $Y_i$ and constant at $f_i(x_i)$ on the other $Y_j$. These two functors form a commutative square: on the selected summand this is tautological, and on each other summand both composites are the same constant functor. Their composite with the inclusion is the identity of $f_i$. Thus $f_i$ is a retract of $\coprod f_j$ and belongs to $W$. \hfill$\square$

This proves a sharper statement than the finite reflection needed in Proposition 9.4. It uses no lifting property and no properness theorem. The empty-source case is the only point where preservation of components is needed; the chosen point supplies the retraction in every other case.

Corollary 9.6 can now be read directly at the level of components, even without this retract argument. Under Loc\,(5), an arrow in $W$ restricts to an arrow in $W$ over each connected component of its target. If every such arrow with connected nonempty target has connected nonempty source, these restrictions each have exactly one source component; therefore the original map induces a bijection on components. The converse follows by applying a component-preserving map to a target having one component. Likewise, the equivalence between (a) and (c) of Proposition 9.11 is formal: a component bijection identifies the coproduct decompositions, and Loc\,(5\,bis) applies to that coproduct. Conversely, criterion (c) implies component preservation and reconstructs both preservation and reflection of arbitrary coproducts by inspecting their component maps. Only the comparison with nondegeneracy in (b) invokes the collapse theorem discussed below.

For a fundamental localizer in the modern, relative sense, coproduct closure and retract closure are available. Hence Loc\,(8) implies Loc\,(5\,bis) in that setting. The reverse implication fails: the class of all functors satisfies decomposition and does not preserve components. Thus the proposed mutual independence does not survive in the modern setting. For the literal weaker assumptions Loc\,(1,2,3) of the present manuscript, the cited modern theorems do not eliminate the distinction between weak and relative localization; we do not claim that they settle that separate implication.

There is also a direct explanation of footnote 57. In any category with sequential colimits, a retract $a$ of an object $b$, with inclusion $i$ and retraction $r$, is the colimit of the constant-object sequence
\[
 b\xrightarrow{ir}b\xrightarrow{ir}b\xrightarrow{ir}\cdots.
\]
Indeed, a compatible cocone is one arrow $h:b\longrightarrow z$ satisfying $h=hir$, and it factors uniquely through $r$ as $(hi)r$. Apply this observation in the arrow category. A class closed under filtered colimits and isomorphism is therefore closed under retracts. Loc\,(7) supplies the extra saturation condition Loc\,(1\,bis); it does not by this argument alone establish strong saturation.
\end{dcomment}

\begin{dcomment}{examples}{p129}{The two degenerate homotopy categories}
Grothendieck's conclusion allows two outcomes because the empty category can retain information after every nonempty category has become indistinguishable. Both outcomes occur, and computing them helps separate a statement about a localization from a statement about its specified class of denominators.

Let $W_{-1}$ consist of all functors whose source and target are both nonempty, together with the identity of the empty category. Equivalently, these are precisely the functors which preserve and reflect emptiness. This class satisfies two-out-of-three and the inverse test by applying the functor
\[
 E:\Cat\longrightarrow[1],\qquad
 E(A)=\begin{cases}0&A=\emptyset,\\1&A\ne\emptyset.\end{cases}
\]
It satisfies the terminal-object axiom. It even satisfies the relative Theorem A: if $A/s\longrightarrow B/s$ preserves emptiness for every $s\in S$ and $B$ is nonempty, choose $b\in B$ and $s=p(b)$. Then $B/s$ contains $(b,1_s)$, so $A/s$ is nonempty and therefore $A$ is nonempty. The reverse nonemptiness implication is automatic for a functor. Thus the comparison itself lies in $W_{-1}$.

The localization at $W_{-1}$ is $[1]$. To check the universal property, let $F:\Cat\longrightarrow\mathcal H$ invert $W_{-1}$. For nonempty $A$, the arrow $A\longrightarrow e$ becomes invertible. For any $u:A\longrightarrow B$ between nonempty categories, the equation $p_Bu=p_A$ forces
\[
 F(u)=F(p_B)^{-1}F(p_A).
\]
Consequently every arrow between nonempty categories has a uniquely prescribed image, and the unique arrow $\emptyset\longrightarrow e$ supplies the remaining arrow of $[1]$. This is exactly the required factorization, with its natural-isomorphism version as well. Localizing at all functors additionally makes $\emptyset\longrightarrow e$ invertible and gives the terminal category.

The class $W_{-1}$ fails component preservation: the fold $e\amalg e\longrightarrow e$ belongs to it. It also fails summand reflection, since
\[
 (\emptyset\longrightarrow e)\amalg(1_e:e\longrightarrow e)
\]
has nonempty source and target, while its first summand is not in $W_{-1}$. This is the precise obstruction behind the concern on ms.~121. A weak equivalence of a disjoint union need not restrict to a weak equivalence merely because both unions are nonempty.

By contrast, for $W_0=\{u\mid\pi_0(u)\text{ is bijective}\}$, the localization is $\mathrm{Set}$. The component functor has a section sending a set to its discrete category, and the canonical functor $A\longrightarrow\pi_0(A)$ belongs to $W_0$. Its universal property identifies the localization with $\mathrm{Set}$. This supplies infinitely many isomorphism classes, including the finite discrete categories of every cardinality. Thus the alternatives of Proposition 9.12 already occur in these elementary examples; the point of its argument is to show that loss of components leaves no intermediate nonempty homotopy theory under the necessary tensor hypothesis.
\end{dcomment}

\begin{dcomment}{exposition}{p133:Donc, la proposition 9.12 résulte du corollaire 9.13}{The precise hypothesis behind the collapse argument}
Grothendieck has found a remarkably economical way to detect nontriviality: two points must remain distinguishable. The argument is correct once product with a category carries weak equivalences to weak equivalences. This is exactly the operation used when an equality of two points is promoted to an equality of two maps with arbitrary source. We give the full argument with that hypothesis visible, rather than silently proving a statement under stronger axioms than those printed.

Assume Loc\,(1,2,3), and in addition
\[
 u\in W\quad\Longrightarrow\quad u\times1_X\in W
 \quad\text{for every small category }X.\tag{T}
\]
Write $L:\Cat\longrightarrow\mathcal H=\Cat[W^{-1}]$. Natural transformations already identify functors in $\mathcal H$ under the weak axioms. To see this directly, the projection $[1]\times X\longrightarrow X$ has comma category $[1]\times(X/x)$ over $x$, which has a terminal object. Loc\,(2,3) therefore puts the projection in $W$. Its two sections have the same image, being inverses to the same isomorphism. A natural transformation $F\Rightarrow G$ is a functor $[1]\times X\longrightarrow Y$, and composition with the sections gives $L(F)=L(G)$. A finite zigzag of natural transformations gives the same conclusion.

Suppose first that $u:I\longrightarrow J$ belongs to $W$ and identifies two different components of $I$. Choose $i_0,i_1$ in those components, with images $j_0,j_1$ in the same component of $J$. Given $f_0,f_1:X\longrightarrow Y$, split $I=I_0\amalg I_1$ with the chosen objects in different summands, and define $h:I\times X\longrightarrow Y$ by $f_0$ on the first summand and $f_1$ on the second. By (T), $L(u\times1_X)$ is invertible. The two inclusions $X\longrightarrow J\times X$ at $j_0,j_1$ are joined by a zigzag of natural transformations, hence equal in $\mathcal H$. Their inverse images through $L(u\times1_X)$ are equal, and composing with $L(h)$ proves $L(f_0)=L(f_1)$.

Suppose instead that $u$ misses a component of $J$. Let $\varepsilon$ be the discrete category $\{0,1\}$. Define $g_0:J\longrightarrow\varepsilon$ to be constantly zero, and $g_1$ to be zero on every component met by $u$ and one on a chosen complementary nonempty union. The equation $g_0u=g_1u$ and invertibility of $L(u)$ imply $L(g_0)=L(g_1)$. Evaluation at an object of the complementary part gives equality of the two arrows $e\longrightarrow\varepsilon$ in $\mathcal H$. By (T), product with $X$ induces an endofunctor of $\mathcal H$, so the two coproduct inclusions $X\longrightarrow X\amalg X$ have equal images. Composing with $(f_0,f_1):X\amalg X\longrightarrow Y$ again gives $L(f_0)=L(f_1)$. Notice the target $Y$: the printed $X$ in the final display is a slip.

It remains to justify why equality of ordinary functors determines all arrows of the localization. Let $\mathcal Q$ have the objects of $\Cat$, one arrow between any two nonempty objects, one arrow from the empty object to every object, and no arrow from a nonempty object to the empty one. Equality of all parallel ordinary functors gives a factorization of $L$ through $\mathcal Q$. If $W$ contains no arrow with empty source and nonempty target, the functor $\Cat\longrightarrow\mathcal Q$ inverts $W$ and has the same universal property as $L$. Thus $\mathcal H\simeq[1]$. Otherwise invert the unique missing type of arrow in $\mathcal Q$; the resulting category is equivalent to $e$ and has the universal property. In either case the full subcategory on the nonempty categories is equivalent to $e$. This argument includes arbitrary zigzags and does not assume a calculus of fractions.

If, on the other hand, every element of $W$ preserves components, $\pi_0$ factors through $\mathcal H$. The discrete finite categories then remain pairwise nonisomorphic, and the two points of $\varepsilon$ remain distinct. These observations prove all six equivalences of Proposition 9.12 under (T). In particular, (c)$\Rightarrow$(c$'$) is immediate: a nonconstant map on a hom-set exhibits two distinct arrows between nonempty objects. The argument through $\mathcal Q$ proves the contrapositive of (c$'$)$\Rightarrow$(c). The editorial bracket on ms.~133 bis interchanges these two roles; the original direction called trivial there was correct.

The additional hypothesis is available for relative fundamental localizers; see \citet[Propositions 2.1.3 and 2.1.9]{Maltsiniotis2020}. The modern classification of the two exceptional localizers is \citet[Proposition 9.3.2]{Cisinski2006}. Under the manuscript's literal weak hypotheses, the displayed proof has a gap, already flagged by the edition on ms.~133 quinquies. We found no proof in the checked sources that removes it. Finally, the conclusion $\mathcal H\simeq e$ should not be substituted for $W=\Fl(\Cat)$ without a saturation argument. The ambiguous parenthesis on ms.~133 ter is read with this distinction maintained.
\end{dcomment}

\begin{dcomment}{exposition}{p134}{What the model structure supplies to the derivator}
Grothendieck is returning to the purpose of the factorization theorems: they should make homotopy Kan extensions available for diagrams of every shape. The final paragraph says that the proof will be given next, but the announced paragraph is absent. The theorem must therefore be read as a collected statement with several recalled ingredients, not as a theorem proved in these pages.

There is now a convenient independent route to its main structural conclusions. For an accessible fundamental localizer in the relative modern sense, \citet[Theorem 5.2.15]{Cisinski2006} constructs a combinatorial model structure on $\Cat$ with weak equivalences $W$. The associated diagram theory is a strong derivator by \citet[Proposition 1.30]{Groth2013}. These are sufficient modern hypotheses, not a verbatim replacement of the manuscript's L-accessibility condition or a proof that weak Loc\,(3) implies relative Loc\,(3\,bis).

Here is how this route accounts for the assertions. The projective and injective model structures on diagram categories have the same objectwise weak equivalences. Projective cofibrant replacements compute left derived Kan extensions; injective fibrant replacements compute right derived Kan extensions. Derived adjunctions give both adjoints to restriction. The pointwise formulas retain the variance in (10.1.1). If $X:I^{\mathrm{op}}\longrightarrow\Cat$ and $u:I\longrightarrow J$, they have the form
\[
 (u_!X)_j\simeq
 \operatorname{hocolim}_{(j/u)^{\mathrm{op}}}X_i,
 \qquad
 (u_*X)_j\simeq
 \operatorname{holim}_{(u/j)^{\mathrm{op}}}X_i.
\]
In the first comma category an object is $j\longrightarrow u(i)$, and in the second it is $u(i)\longrightarrow j$. The opposite categories occur because the coefficient diagrams are contravariant. They should not be omitted by importing a formula written for covariant diagrams.

The model of a diagram also explains evaluation conservativity. Represent a morphism of the diagram homotopy category by a map from a cofibrant replacement to a fibrant replacement. If every evaluated morphism is invertible in the underlying homotopy category, every component of this representing map is a weak equivalence. Hence the map is an objectwise weak equivalence and is invertible in the diagram homotopy category. This is the content needed for Der~2; it is not a formal consequence of strong saturation of $W$ alone. Compare the commentary on Chapter~VI, manuscript p.~130.

Finite decomposition requires much less. A diagram on $I\amalg J$ is a pair of diagrams, and localization at the corresponding product class gives the product of the two localizations when identities belong to that class. This is the finite-product localization argument already discussed in the commentary on Chapter~X, manuscript p.~5 (``Finite products and sums already descend''). It does not require the coproduct axiom Loc\,(5) for functors in $\Cat$: a coproduct of \emph{index categories} and a coproduct of \emph{coefficient objects} are different operations.

Finally, shifts $I\mapsto\mathbf D(S\times I)$ are derivators, by \citet[Theorem 1.25]{Groth2013}. Pointed versions therefore have the coherent fibre, cofibre, loop and suspension constructions needed for the sequences mentioned here. No assertion of stability is involved: for spaces, suspension and looping are not inverse equivalences, and cocartesian squares need not be cartesian. The older exactness axioms should be read in their pointed, unstable sense. The limitation to certain finite shapes for $u_*$ is superseded under the modern model hypotheses; footnote 58 itself records that Grothendieck later removed it.
\end{dcomment}

\begin{dcomment}{examples}{p135}{Constant objects retain coherent maps}
The detection statement in 10.1(3a) becomes natural when we ask what the morphisms between constant diagrams actually are. The full subcategory of constant objects has all the morphisms inherited from the homotopy category of diagrams. These morphisms remember the indexing homotopy type, even though the displayed objects are constant. This is the same issue introduced in the commentary on Chapter~I, manuscript pp.~61--62 (``Equivalence measured by constant coefficients'' and ``How much constant coefficients can detect''); the following computation makes its role in the present theorem explicit.

For the ordinary homotopy theory of spaces, write $BS=|N(S)|$ and let $K,L$ be spaces. The derived space of maps between their constant diagrams on $S$ is the homotopy limit of the constant mapping-space diagram, hence
\[
 \operatorname{Map}_{\operatorname{Fun}(S^{\mathrm{op}},\mathcal S)}
       (\operatorname{const}K,\operatorname{const}L)
 \simeq \operatorname{Map}(BS,\operatorname{Map}(K,L)).
\]
On connected components this gives
\[
 \Hom_{\mathbf D(S)}(K_S,L_S)\cong[BS\times K,L].\tag{*}
\]
One can obtain the same formula by the adjunction between homotopy colimit and the constant diagram functor: the homotopy colimit of the constant diagram $K$ is $BS\times K$. This derivation also shows that restriction along $f:S'\longrightarrow S$ corresponds to precomposition with $Bf\times1_K$.

If $Bf$ is a weak equivalence, every such map of hom-sets is a bijection. Conversely, take $K=*$ and let $L$ vary. Full faithfulness on constant objects says that
\[
 [BS,L]\longrightarrow[BS',L]
\]
is a bijection for every space $L$. The contravariant Yoneda lemma in the homotopy category then says that $Bf$ is invertible there, hence is a weak equivalence. Essential surjectivity on the full constant subcategories is automatic: the constant object with value $L$ on $S'$ is the restriction of the constant object with value $L$ on $S$. Thus full faithfulness and equivalence give the same test here.

For a concrete calculation let $S=B\mathbb Z$, the one-object category of the group $\mathbb Z$, and let $K=*$, $L=S^1$. Its classifying space is $S^1$. Formula $(*)$ gives
\[
 \Hom_{\mathbf D(B\mathbb Z)}(*_{B\mathbb Z},(S^1)_{B\mathbb Z})
 \cong[S^1,S^1]\cong\mathbb Z.
\]
The integer is the degree of the corresponding map. On restriction along the object $e\longrightarrow B\mathbb Z$, all these morphisms map to the unique element of $[*,S^1]$. The distinction between degrees has been forgotten. Consequently this restriction is not fully faithful on constant objects, exactly as the theorem predicts. Looking only at strict natural transformations between the original constant diagrams would miss these derived morphisms.

For a general $W$, the analogous proof uses the homotopy-colimit formula for constant families, product stability, invariance under passage to opposite categories, and strong saturation. These ingredients are available in the relative accessible setting of the preceding block. The smaller hypothesis list in footnote 60 should be read with the earlier construction it invokes, not as a new proof of those ingredients under weak Loc\,(3). In particular, strong saturation is what changes ``the comparison is invertible in $\Hot_W$'' into ``$f\in W$''.
\end{dcomment}

\begin{dcomment}{development}{p136}{A finite monodromy test for the displayed hypotheses}
The extra hypothesis in 10.1(3b) is meant to control locally constant families, whose fibres may be transported nontrivially around loops. There is a consequential numbering slip: the printed Loc\,(5\,b) is described as the fibration axiom, but in the new list Loc\,(5\,b) means reflection of coproducts. The intended reference is Loc\,(6\,b), corresponding to W\,(4\,b) in chapter XII. The distinction can be tested by a finite group.

\textbf{Failure under the literal reduced list.}\ours
Let $W_{\mathbb Q}$ be the functors inducing isomorphisms on rational homology of nerves. It satisfies strong saturation, the terminal-object axiom, relative Theorem A and coproduct reflection. Nevertheless there is $f\in W_{\mathbb Q}$ for which restriction on locally constant objects is not essentially surjective.

\emph{Proof.} Strong saturation follows directly from the definition: rational homology factors through the localization, so an arrow made invertible by localization must induce an isomorphism on rational homology. The terminal-object assertion follows from contractibility of the nerve. Rational homology of a coproduct is the direct sum of the homologies of its summands. A block-diagonal direct sum map is invertible precisely when each component map is invertible, which proves coproduct preservation and reflection. That $W_{\mathbb Q}$ is a relative fundamental localizer follows from \citet[Corollary 9.4.3 and \S9.4.4]{Cisinski2006}; equivalently, it is the nerve pullback of the homological localization of simplicial sets.

Choose a nontrivial finite group $G$, for example $C_2$. The map $p:BG\longrightarrow e$ belongs to $W_{\mathbb Q}$. Indeed, the augmentation $\mathbb Q[G]\longrightarrow\mathbb Q$ splits as a map of $\mathbb Q[G]$-modules by
\[
 1\longmapsto\frac1{|G|}\sum_{g\in G}g.
\]
Thus the trivial module $\mathbb Q$ is projective, and
\[
 H_n(BG;\mathbb Q)=
 \operatorname{Tor}^{\mathbb Q[G]}_n(\mathbb Q,\mathbb Q)=0\quad(n>0),
 \qquad H_0(BG;\mathbb Q)=\mathbb Q.
\]
Let $A:(BG)^{\mathrm{op}}\longrightarrow\Cat$ be the discrete category on the regular right $G$-set, with its action. Every structure map is an isomorphism, so $A$ is locally constant.

The component functor sends every $W_{\mathbb Q}$-map to a bijection. To verify this last point without a finiteness assumption, $H_0(NC;\mathbb Q)$ is the vector space with basis $\pi_0(C)$, and a map induced by a function on those bases is invertible only when that function is bijective: two identified basis vectors give a nonzero kernel element, while an omitted basis vector cannot lie in the image. Applied objectwise, $\pi_0$ therefore descends to a functor
\[
 \Cat^{(BG)^{\mathrm{op}}}[W_{\mathbb Q}(BG)^{-1}]
 \longrightarrow \mathrm{Set}^{(BG)^{\mathrm{op}}}.
\]
An object in the image of $p^*$ has trivial $G$-action after this functor is applied. The regular $G$-set is nonempty and has no point fixed by all of $G$, so it is not equivariantly isomorphic to any nonempty trivial $G$-set. Hence $A$ cannot be isomorphic to a constant object in the localized diagram category. This proves the failure of essential surjectivity. \hfill$\square$

The example also satisfies the opening excision and size hypotheses of 10.1. For Loc\,(4), a simplex in the nerve of a category covered by two open subcategories lies in one of the two subcategories: for sieves, its last vertex determines this, and for the opposite convention its first vertex does. Consequently
\[
 C_*(NX,NX_2;\mathbb Q)
 \cong C_*(NX_1,NX_0;\mathbb Q).
\]
The two inclusions $X_0\longrightarrow X_1$ and $X_2\longrightarrow X$ are rational homology equivalences precisely when these isomorphic relative chain complexes are acyclic. This proves both directions of Loc\,(4). Unlike the corresponding assertion for ordinary weak homotopy equivalences, this reverse implication is valid for the homological invariant used here.

For completeness, L-accessibility can be checked in the manuscript's own sense. Fix an infinite cardinal $\pi$ and a rational homology equivalence $f:X\longrightarrow Y$. Start with a subarrow $f_0:X_0\longrightarrow Y_0$ having at most $\pi$ arrows in total. The rational mapping cone $C_*(Nf_0;\mathbb Q)$ has at most $\pi$ elements. For every cycle in it choose a bounding chain in the acyclic cone $C_*(Nf;\mathbb Q)$. Each bounding chain uses finitely many simplices. Adjoin all their objects and arrows, their images under $f$, and all necessary composites. This gives a subarrow $f_1$ still of size at most $\pi$. Repeat countably often. In the union $\bar f$, every cycle has finite support and therefore occurs at one stage; its chosen bounding chain occurs at the next. Thus $\bar f\in W_{\mathbb Q}$ and has size at most $\pi$.

It follows that the $W_{\mathbb Q}$-subarrows of size at most $\pi$ are cofinal among all subarrows of that size and form a filtered family whose union is $f$. Rational chain complexes of nerves commute with filtered colimits, and filtered colimits of vector spaces preserve exactness, so $W_{\mathbb Q}$ is closed under all filtered colimits. There is a set of representatives of the arrows of size at most $\pi$. They generate $W_{\mathbb Q}$ under this closure. This is precisely the condition displayed in Loc\,(7\,bis). It also gives the stronger bounded-subobject criterion used in the second Proposition 5.7: categories and functors of size at most $\pi$ have presentations with at most $\pi$ generators and relations, hence maps out of them factor through a stage of a $\pi^+$-filtered diagram, with equalities detected at a later stage. Thus the comparison with manuscript L-accessibility is explicit here; it is not inferred from the modern terminology.

The example therefore refutes the assertion with current Loc\,(5\,b) read literally, even together with the opening assumptions. It does not refute the corrected assertion with the intended fibration axiom. The exact issue is that a test which preserves components of disjoint unions need not eliminate monodromy.

The regular action makes the missing fibration hypothesis visible as well. Its translation groupoid $EG$ has one arrow between any two objects and is contractible. The projection $EG\longrightarrow BG$ has discrete fibre $G$. It is an old-style $W$-fibration: on a contractible base the discrete covering is trivializable, and a map between two such bases induces weak equivalences on the corresponding sheets. Compare it over $BG$ with the identity fibration $BG\longrightarrow BG$. The map of total categories $EG\longrightarrow BG$ is a rational homology equivalence, while its fibre map $G\longrightarrow e$ is not. This directly violates Loc\,(6\,b). Pullback along $e\longrightarrow BG$ also gives $G\longrightarrow EG$, which is not a rational homology equivalence because $|G|>1$, illustrating failure of the base-change condition Loc\,(6\,a). We checked the manuscript's earlier fibration variants at \citet[chapter XII, ms.~53--56]{Grothendieck1990XIIa} and the rational-localizer example in Cisinski's \S9.4; we did not find this precise diagnostic of the renumbered assertion stated there. The proof is offered as a check of the manuscript, not as a new claim about homological localization.
\end{dcomment}

\begin{dcomment}{development}{p136}{Why a nonempty shift still detects the base test}
Footnote 61 asks that the shifted assertion be checked with care. One half is formal and explains exactly why the exclusion of the empty category is necessary. The other half needs a theorem about coherent diagrams with the shifted coefficient category; checking each value separately does not by itself prove full faithfulness of the whole diagram functor.

\textbf{Reflection by a nonempty shift.}\ours
Let $\mathbf D$ be a prederivator, let $I$ have an object $i$, and put $\mathbf D^I(S)=\mathbf D(I\times S)$. If restriction along $f:S'\longrightarrow S$ is fully faithful on the full subcategories of constant objects for $\mathbf D^I$, then it is fully faithful on the corresponding full subcategories for $\mathbf D$.

\emph{Proof.} The projection $I\times S\longrightarrow S$ and the inclusion $s\mapsto(i,s)$ induce functors
\[
 c_S:\mathbf D(S)\longrightarrow\mathbf D(I\times S),
 \qquad e_S:\mathbf D(I\times S)\longrightarrow\mathbf D(S)
\]
with $e_Sc_S=1$, up to the canonical coherence isomorphism if pseudofunctors are used. These functors commute with restriction along $f$. If $A,B\in\mathbf D(e)$, the objects $c_S(A_S),c_S(B_S)$ are constant in the $S$-direction, with values $c_e(A),c_e(B)\in\mathbf D(I)$. Therefore the map
\[
 \Hom_{\mathbf D(S)}(A_S,B_S)
 \longrightarrow\Hom_{\mathbf D(S')}(A_{S'},B_{S'})
\]
is a retract, in the arrow category of sets, of the corresponding map between the hom-sets of the shifted constant objects. The latter is bijective by hypothesis. A retract of a bijection is bijective: if $q$ is the large bijection, its inverse on the retract is obtained by inclusion, then $q^{-1}$, then retraction. The two commutative squares and the retraction identities verify both inverse equations. This proves full faithfulness. \hfill$\square$

Every constant object on $S'$ is the restriction of the constant object with the same coefficient on $S$, so the test is also reflection of equivalence between the full constant subcategories. The argument uses an object of $I$, not an initial or terminal object, and it does not require $I$ to be connected. If $I=\emptyset$ and $\mathbf D$ satisfies the usual empty-index axiom, every category $\mathbf D^I(S)$ is terminal; every $f$ then passes the test, so detection of $W$ is impossible.

For the ordinary localizer $W_\infty$, the forward statement and its locally constant refinement hold for every shift: apply \citet[Theorem 1.3]{Cisinski2009LocallyConstant} to a model category of $I$-diagrams. This coefficient category retains coherent diagrams; it is not merely the ordinary functor category with values in $\Hot$. The theorem just cited concerns nerve weak equivalences. It is not a theorem that every element of an arbitrary larger localizer preserves all locally constant coefficients; the preceding finite-group example explains why such a reformulation would be false. We checked this formal shifted argument against \citet[Theorem 1.25]{Groth2013} and the derivator discussion in Cisinski's \S3, and supply the retract proof to settle the specific reflection point left implicit here.
\end{dcomment}

\begin{dcomment}{later}{p137}{The later theorem and the scope of the unfinished proof}
The principal construction announced here has a full modern realization through model categories and their diagram homotopy theories. Locally constant coefficients likewise have a precise invariant formulation: \citet[Theorem 1.3 and Corollaries 1.28 and 3.8]{Cisinski2009LocallyConstant} proves invariance under nerve weak equivalences, and conversely characterizes those equivalences by testing all model-category or derivator coefficients. In the published paper these are the numbers; the shorter preprint has different numbering. In higher-category language the underlying object is the functor category out of the groupoid completion of the indexing category, which is the perspective described in that paper's introduction.

The unresolved matters in these pages are therefore specific. The missing next paragraph does not provide a proof of 10.1 under each of its reduced hypothesis lists. The weak-versus-relative localization issue remains visible in 9.12 and in the claimed construction of $u_!$. The fibration reference in 10.1(3b) must be read after correcting its old numbering. None of these points is a reason to withhold the established constructions under their verified hypotheses, and none is repaired simply by citing a theorem about the ordinary localizer while leaving an arbitrary $W$ in the statement.
\end{dcomment}
